% !TEX program = xelatex
\documentclass[11pt,a4paper]{article}
\usepackage{fontspec}
\setmainfont{texgyrepagella}[Extension=.otf,UprightFont=*-regular,BoldFont=*-bold,ItalicFont=*-italic,BoldItalicFont=*-bolditalic]
\usepackage{newunicodechar}
\newunicodechar{≈}{\ensuremath{\approx}}
\newunicodechar{→}{\ensuremath{\rightarrow}}
\newunicodechar{≥}{\ensuremath{\geq}}
\newunicodechar{≤}{\ensuremath{\leq}}
\newunicodechar{−}{\ensuremath{-}}
\newunicodechar{×}{\ensuremath{\times}}
\usepackage[italian]{babel}
\usepackage[a4paper,margin=2.5cm]{geometry}
\usepackage{microtype}
\usepackage{booktabs,xltabular,array}
\usepackage{enumitem}
\setlist{itemsep=2pt,topsep=4pt}
\usepackage{xcolor}
\definecolor{accent}{HTML}{3B6A4C}
\usepackage{titlesec}
\titleformat{\section}{\large\bfseries\color{accent}}{\thesection.}{0.6em}{}
\titleformat{\subsection}{\normalsize\bfseries}{\thesubsection}{0.6em}{}
\usepackage{fancyhdr}
\pagestyle{fancy}\fancyhf{}
\fancyhead[L]{\small\textsc{NexaCore · Piano integrato}}
\fancyhead[R]{\small Gruppo 2}
\fancyfoot[C]{\small\thepage}
\renewcommand{\headrulewidth}{0.4pt}
\usepackage[hidelinks]{hyperref}
\setlength{\parindent}{0pt}
\setlength{\parskip}{5pt}

\title{\textbf{NexaCore S.p.A.}\\[4pt]\Large Piano integrato di governance del cambiamento e del rischio}
\author{Paolo Gallignani \quad Marco Aiello \quad Alex Domeniconi\\
Michele Guidetti \quad Gloria Bleve \quad Pietro Chiodi\\[10pt]
\small Gruppo 2 · Challenge Cyber \& Tech Risk KPMG\\
\small Corso di Comportamento Organizzativo, Università di Bologna}
\date{16 novembre 2026}

\begin{document}
\maketitle
\thispagestyle{empty}

\begin{abstract}
\noindent NexaCore (850 dipendenti, 6 sedi in 5 paesi) ha già approvato 50 M€ per un ERP in cloud, l'intelligenza artificiale generativa e l'automazione. Il nostro punto di partenza è che la tecnologia è decisa e il rischio sta nelle persone. Seguendo la ricerca-azione, il piano prima conosce le persone con un trivia su autoefficacia, autodeterminazione ed equità percepita, restituisce loro il profilo e poi agisce: strutture temporanee di governance con KPMG in ruolo consultivo, un pilota ERP su una sola linea di prodotto con due go/no-go, AI generativa per tutti con un contratto enterprise, automazioni dei Champions dentro una classificazione del rischio, formazione e recognition tra pari. Il programma dura 24 mesi, impegna circa 20,4 M€ dei 50 approvati, sbloccati a tranche, e libera a regime circa 70.000 ore l'anno. La valutazione combina KPI di continuità e di rischio con misure di commitment al cambiamento e di benessere.
\end{abstract}

\newpage
\tableofcontents
\newpage


\section*{Come leggere questo documento}

Questo è il documento completo del piano. La presentazione (5 slide) ne è la sintesi; qui c'è tutto ciò che la presentazione non ha spazio per spiegare: le ragioni delle scelte, i numeri, le assunzioni e i loro limiti.

È scritto per essere letto da una persona oppure interrogato con un assistente di intelligenza artificiale. Per questo:
\begin{itemize}
\item ogni sezione è numerata e si può citare (es. ``\hyperref[sec:9.4]{\S9.4}'');
\item i termini usati sono sempre gli stessi e sono definiti nel glossario (\hyperref[sec:21]{\S21});
\item \textbf{NexaCore è un'azienda fittizia}: i dati del caso sono quelli forniti da KPMG; tutto il resto (numeri di persone per ruolo, costi, soglie) sono \textbf{stime e assunzioni del gruppo}, dichiarate come tali. Nessun dato è stato raccolto in un'azienda reale;
\item le assunzioni sono marcate con lo stato: \textbf{assunzione} (scelta del gruppo per poter progettare), \textbf{proposta} (da confermare con la governance di NexaCore), \textbf{da verificare} (dipende da norme, contratti o dati non disponibili);
\item i mesi si contano dall'avvio del programma: M0 è la situazione di partenza, M1 il primo mese di lavoro, M24 l'ultimo;
\item i costi sono in milioni di euro (M€), al valore centrale quando esiste una forchetta.
\end{itemize}

Chi cerca una risposta rapida può partire dalla sintesi (\hyperref[sec:1]{\S1}), dalla matrice dei timori del Board (\hyperref[sec:17]{\S17}) o dalle domande frequenti (\hyperref[sec:20]{\S20}).

\section{Sintesi}\label{sec:1}

\textbf{Il problema.} NexaCore (850 dipendenti, 6 sedi in 5 paesi, 8 reparti) ha già approvato 50 M€ per tre tecnologie: un ERP in cloud, l'intelligenza artificiale generativa (GenAI) e l'automazione dei processi. Il Consiglio di Amministrazione teme cinque rischi: resistenza al cambiamento, perdita di conoscenze chiave, errori operativi nella transizione, nuovi rischi cyber e uso non governato dell'AI. Il caso chiede un modello integrato di governance del cambiamento e del rischio che protegga continuità operativa, controllo dei rischi tecnologici e valorizzazione delle persone.

\textbf{La nostra tesi.} La tecnologia è decisa; il rischio sta nelle persone. Quattro dei cinque timori del Board dipendono da come le persone accolgono, usano e trasmettono le nuove tecnologie, e anche il quinto (cyber) si gioca in buona parte sui comportamenti. Per questo il piano parte dal conoscere le persone e solo dopo agisce.

\textbf{Il metodo.} Seguiamo la \textbf{ricerca-azione}, l'approccio moderno al cambiamento organizzativo: prima si capisce il problema insieme a chi lo vive (diagnosi, analisi, feedback), poi si agisce, si misura e si corregge (azione, valutazione).

\textbf{Le scelte principali.}
\begin{enumerate}
\item \textbf{Un trivia obbligatorio e nominativo} (12 domande su autoefficacia, autodeterminazione ed equità percepita) colloca ogni persona in uno di 8 profili e quindi tra \textbf{Champions} (≈20\%), \textbf{Fence sitters} (≈70\%) e \textbf{Cinici} (≈10\%). Serve ad assegnare i ruoli nel cambiamento, mai a valutare le persone. Ognuno riceve il proprio profilo e può contestarlo.
\item \textbf{La struttura gerarchica non cambia.} Si aggiungono strutture temporanee: un \textbf{team di governance} che decide (vertice aziendale, CISO, 2 Champions, consulenti KPMG senza voto), un \textbf{Change \& Risk Office} che coordina, ≈40 \textbf{referenti} nelle funzioni.
\item \textbf{L'ERP parte da una sola linea di prodotto} (≈65 persone, M2-M10), il team con più Champions, con KPMG che configura e poi esce. Due \textbf{go/no-go}: a M6 per accendere il sistema, a M10 per estenderlo. Poi \textbf{4 ondate} sulle altre linee (M11-M22), mai più di due linee in transizione insieme.
\item \textbf{L'AI generativa arriva a tutti} con un contratto enterprise (i dati restano di NexaCore e non addestrano i modelli), tre profili d'accesso e formazione di 4, 8 o 24 ore.
\item \textbf{32 Champions fuori dal pilota automatizzano} le attività ripetitive del proprio lavoro, dentro una \textbf{classificazione del rischio} decisa dalla governance, con un \textbf{registro} e un \textbf{sostituto} per ogni automazione.
\item \textbf{Cinque comportamenti osservabili} (decido con i dati, uso l'AI entro le regole, documento ciò che so, segnalo errori e incidenti, aiuto un collega) sono il cuore della cultura e si rinforzano con una \textbf{piattaforma di recognition} tra colleghi.
\item \textbf{Valutazione continua}: KPI di continuità e rischio, commitment al cambiamento (Herscovitch e Meyer) e benessere (WHO-5, HSE, Copenhagen Burnout Inventory), con misure a M1, M10, M16 e M23-M24.
\end{enumerate}

\textbf{I numeri chiave.}

{\small\renewcommand{\arraystretch}{1.25}
\begin{xltabular}{\linewidth}{@{}>{\raggedright\arraybackslash\hsize=0.525\hsize\hspace{0pt}}X>{\raggedright\arraybackslash\hsize=1.475\hsize\hspace{0pt}}X@{}}
\toprule
\textbf{Voce} & \textbf{Valore} \\
\midrule
\endhead
Durata del programma & 24 mesi: 22 di lavoro + 2 di chiusura \\
Budget pianificato & ≈20,4 M€ su 50 approvati (41\%), riserva del 20\% compresa \\
Tranche 1 (M1-M10) & ≈7 M€, all'avvio \\
Tranche 2 (M11-M24) & ≈13 M€, solo se il go/no-go di M10 è superato \\
Non impegnato & ≈30 M€, restano a disposizione del CdA \\
Benefici a regime & ≈70.000 ore l'anno liberate (≈45 persone a tempo pieno, ≈2,1 M€ l'anno), reinvestite, nessun taglio di posti \\
\bottomrule
\end{xltabular}}

\textbf{Rispondiamo ai tre output richiesti dal caso.}
\begin{itemize}
\item \emph{Modello di governance:} \hyperref[sec:8]{\S8} (strutture, ruoli, decisioni), \hyperref[sec:14]{\S14} (go/no-go e tranche).
\item \emph{Azioni pratiche:} \hyperref[sec:9]{\S9}-§13.
\item \emph{KPI per misurare l'impatto:} \hyperref[sec:16]{\S16} e matrice in §17.
\end{itemize}

\section{Il caso e il nostro approccio}\label{sec:2}

\subsection{I dati del caso}\label{sec:2.1}
\begin{itemize}
\item \textbf{NexaCore S.p.A.}: multinazionale di servizi tecnologici e soluzioni digitali per le imprese.
\item \textbf{850 dipendenti}: 300 sotto i 30 anni, 370 tra 30 e 50, 180 sopra i 50.
\item \textbf{6 sedi in 5 paesi}: Italia, Germania, Francia, Regno Unito, Spagna.
\item \textbf{8 reparti}: Operations, Finance \& Administration, Sales, Marketing, Human Resources, IT \& Digital Transformation, Procurement, Customer Service.
\item \textbf{Modello di lavoro}: 70\% in presenza, 30\% da remoto.
\item \textbf{Investimento approvato}: 50 M€ per ERP in cloud, GenAI e automazione.
\item \textbf{Timori del Board} (elenco aperto nel caso): resistenza al cambiamento, perdita di conoscenze chiave, errori operativi durante la transizione, nuovi rischi cyber, uso non governato dell'AI.
\item \textbf{Obiettivo}: un modello integrato di governance del cambiamento e del rischio che garantisca continuità operativa, controllo dei rischi tecnologici e valorizzazione delle persone, integrando gestione del cambiamento, governance e risk management, knowledge management e sviluppo delle competenze, con indicatori di impatto.
\item \textbf{Criteri di valutazione}: coerenza e completezza, creatività, fattibilità, misurabilità.
\end{itemize}

\subsection{Perché la ricerca-azione}\label{sec:2.2}
L'approccio classico (Lewin, Kotter) vede il cambiamento come una sequenza pianificata: creare la consapevolezza, accompagnare la transizione, consolidare. L'approccio moderno riconosce che in contesti complessi e ambigui il problema va prima capito insieme a chi lo vive. La \textbf{ricerca-azione} fa proprio questo: raccoglie dati in modo sistematico, li restituisce alle persone, e solo dopo definisce gli interventi, che poi misura e corregge.

Il suo ciclo — \textbf{diagnosi}, \textbf{analisi}, \textbf{feedback}, \textbf{azione}, \textbf{valutazione} — dà la struttura a questo documento e resta visibile sulla roadmap: diagnosi e analisi (M1), feedback (M2), azione (M2-M22), valutazione (M10, M16, M23-M24). Non è una sequenza che si percorre una volta sola: a M10 la valutazione di fine pilota riapre il ciclo.

\subsection{Le cinque fasi del piano}\label{sec:2.3}

{\small\renewcommand{\arraystretch}{1.25}
\begin{xltabular}{\linewidth}{@{}>{\raggedright\arraybackslash\hsize=0.455\hsize\hspace{0pt}}X>{\raggedright\arraybackslash\hsize=2.256\hsize\hspace{0pt}}X>{\raggedright\arraybackslash\hsize=0.288\hsize\hspace{0pt}}X@{}}
\toprule
\textbf{Fase} & \textbf{Domanda a cui risponde} & \textbf{Dove} \\
\midrule
\endhead
\textbf{Diagnosi} & Com'è fatta l'organizzazione e chi sono le persone? & \hyperref[sec:4]{\S4}, \hyperref[sec:5]{\S5} \\
\textbf{Analisi} & Dove nascono i problemi e la resistenza? & \hyperref[sec:6]{\S6} \\
\textbf{Feedback} & Le persone si riconoscono nel quadro che ne abbiamo fatto? & \hyperref[sec:7]{\S7} \\
\textbf{Azione} & Che cosa facciamo, chi lo fa, quando? & \hyperref[sec:8]{\S8}-\hyperref[sec:15]{\S15} \\
\textbf{Valutazione} & Funziona? Come lo vivono le persone? & \hyperref[sec:16]{\S16} \\
\bottomrule
\end{xltabular}}

Il ciclo è \textbf{ricorsivo}: a fine pilota (M10) si ripetono il trivia e le misure, si confrontano con i comportamenti osservati e si correggono ruoli e interventi prima di estendere.

\section{Assunzioni di partenza}\label{sec:3}

Il caso lascia aperti molti dettagli. Per progettare abbiamo fissato poche assunzioni, tutte dichiarate.

{\small\renewcommand{\arraystretch}{1.25}
\begin{xltabular}{\linewidth}{@{}>{\raggedright\arraybackslash\hsize=0.173\hsize\hspace{0pt}}X>{\raggedright\arraybackslash\hsize=2.410\hsize\hspace{0pt}}X>{\raggedright\arraybackslash\hsize=0.417\hsize\hspace{0pt}}X@{}}
\toprule
\textbf{\#} & \textbf{Assunzione} & \textbf{Stato} \\
\midrule
\endhead
A1 & Il software ERP è già scelto e rientra nei 50 M€; KPMG è già ingaggiata per implementarlo e non è il revisore dei conti di NexaCore. (Se il software non fosse scelto, servirebbero 2-3 mesi di selezione con un consulente diverso da chi implementa.) & assunzione \\
A2 & La struttura gerarchica è funzionale: CEO e 8 reparti. I responsabili più probabili per un'azienda di queste dimensioni sono CFO, COO, CIO, direttore HR e, dentro IT, un CISO. Non aggiungiamo altre figure. & assunzione \\
A3 & Accanto ai reparti esistono già \textbf{team di prodotto interfunzionali}, uno per linea di prodotto, che attraversano i reparti dall'ordine all'incasso. Sono le uniche strutture stabili del progetto. & assunzione \\
A4 & Le linee di prodotto sono ≈5: una fa da pilota, le altre 4 passano all'ERP in 4 ondate. I clienti appartengono a 4 settori (manifatturiero, energia, logistica, servizi finanziari). & assunzione (vedi \hyperref[sec:19]{\S19}) \\
A5 & Oggi non esiste uno strumento di AI aziendale; alcuni dipendenti usano strumenti personali e l'azienda ne scoraggia l'uso. & assunzione \\
A6 & I 50 M€ si sbloccano \textbf{a tranche}, legate ai go/no-go. & proposta \\
A7 & Le competenze interne bastano per la maggior parte dei ruoli: poche assunzioni (≈8 stabili), sempre dopo aver provato la mobilità interna. & assunzione \\
A8 & Ripartizione iniziale attesa 20/70/10 tra Champions, Fence sitters e Cinici (≈170 / ≈580 / ≈100), da sostituire con i dati del trivia. Con un trivia nominativo ci aspettiamo più Champions sulla carta che nei fatti. & assunzione \\
A9 & Gli utenti dell'ERP a regime sono ≈500 (chi lavora nei processi amministrativi, commerciali e logistici). & assunzione \\
A10 & Costo medio del lavoro ≈30 €/ora, ≈48.000 € l'anno (Eurostat, Italia 2023: 29,8 €/ora). La media pesata dei 5 paesi è probabilmente più alta. & assunzione, da verificare \\
A11 & Le ore liberate dalla tecnologia non diventano tagli di personale: si reinvestono in attività di maggior valore e nella mobilità verso i ruoli nuovi. & proposta \\
\bottomrule
\end{xltabular}}

\section{Diagnosi: l'organizzazione e i timori del Board}\label{sec:4}

\subsection{Struttura hard: funzionale, e resta così}\label{sec:4.1}
L'organizzazione per reparti favorisce specializzazione, ruoli chiari e controllo. Il suo limite sono i \textbf{compartimenti} (silos) e il coordinamento difficile tra funzioni: proprio quello che l'ERP mette alla prova, perché unisce in un solo sistema dati e processi che oggi ogni reparto gestisce a modo suo.

\textbf{Scelta:} non riorganizziamo. Cambiare la gerarchia nel mezzo di tre transizioni tecnologiche aumenterebbe l'incertezza e il rischio operativo. Il coordinamento orizzontale lo garantiscono i team di prodotto già esistenti e alcune strutture temporanee (\hyperref[sec:8]{\S8}).

\subsection{Struttura soft: i team di prodotto interfunzionali}\label{sec:4.2}
La struttura soft è fatta di relazioni, flussi di comunicazione e collaborazione oltre l'organigramma. In NexaCore, con 6 sedi, 5 paesi e il 30\% di lavoro da remoto, la capacità di condividere informazioni e conoscenze è critica. Assumiamo (A3) che esistano già team di prodotto che riuniscono persone di reparti diversi attorno alla stessa linea. Il piano li usa come \textbf{unità del cambiamento}: il pilota ERP si fa dentro uno di essi, e le ondate seguono le linee.

\subsection{I timori del Board e le loro radici}\label{sec:4.3}

{\small\renewcommand{\arraystretch}{1.25}
\begin{xltabular}{\linewidth}{@{}>{\raggedright\arraybackslash\hsize=0.360\hsize\hspace{0pt}}X>{\raggedright\arraybackslash\hsize=1.640\hsize\hspace{0pt}}X@{}}
\toprule
\textbf{Timore} & \textbf{Da dove nasce} \\
\midrule
\endhead
\textbf{Resistenza al cambiamento} & Abitudini consolidate, bisogno di sicurezza, paura per il ruolo; inerzia di struttura e di gruppo. \\
\textbf{Perdita di conoscenze chiave} & Il sapere di processo è spesso tacito, in testa agli esperti. Chi teme di diventare sostituibile è meno incline a condividerlo. \\
\textbf{Errori operativi nella transizione} & Saper usare un sistema non significa averlo assimilato: restano file e procedure parallele; dati incerti tra reparti. \\
\textbf{Nuovi rischi cyber} & Cloud, integrazioni e AI ampliano la superficie d'attacco e il numero di accessi da governare. \\
\textbf{AI non governata} & Se lo strumento ufficiale manca o è scomodo, le persone usano strumenti personali (\emph{shadow AI}), con dati aziendali fuori controllo. \\
\bottomrule
\end{xltabular}}

\subsection{Strumenti della diagnosi}\label{sec:4.4}
\begin{enumerate}
\item \textbf{Trivia} per tutti (\hyperref[sec:5]{\S5}): profila le persone su tre dimensioni motivazionali.
\item \textbf{Interviste} a dipendenti e responsabili, prima e dopo il trivia: raccolgono percezioni e difficoltà e verificano che i profili corrispondano all'esperienza. Si ripetono durante e dopo il pilota.
\item \textbf{Benchmark} su casi di trasformazioni analoghe: Robey et al. (2002) sulle barriere di conoscenza nell'ERP, SACE, HolidayCo, Lidl, Accenture, IBM, TCS, Qantas.
\end{enumerate}

Le interviste e il trivia sono \textbf{attività previste}: questo documento non contiene risultati raccolti in NexaCore.

\section{Diagnosi: il trivia e le otto personas}\label{sec:5}

\subsection{Perché tre dimensioni}\label{sec:5.1}
Fiducia nelle proprie capacità, ragioni per impegnarsi ed equità percepita possono divergere, e chiedono interventi diversi: chi vede il valore del cambiamento può aver bisogno di pratica; chi si sente capace può frenare se giudica sbilanciato il carico.

{\small\renewcommand{\arraystretch}{1.25}
\begin{xltabular}{\linewidth}{@{}>{\raggedright\arraybackslash\hsize=0.823\hsize\hspace{0pt}}X>{\raggedright\arraybackslash\hsize=0.385\hsize\hspace{0pt}}X>{\raggedright\arraybackslash\hsize=1.558\hsize\hspace{0pt}}X>{\raggedright\arraybackslash\hsize=1.234\hsize\hspace{0pt}}X@{}}
\toprule
\textbf{Dimensione} & \textbf{Teoria} & \textbf{Domanda di fondo} & \textbf{Che cosa ci dice per l'azione} \\
\midrule
\endhead
\textbf{A · Autoefficacia} & Bandura & Mi sento capace di usare il nuovo strumento? & Quanto supporto e quanta pratica offrire. \\
\textbf{D · Autodeterminazione} & Deci e Ryan & Cambio per interesse o valore (motivazione autonoma), per premi o pressioni (motivazione controllata), o per nulla (amotivazione)? & Quanta autonomia lasciare e quale senso costruire. \\
\textbf{E · Equità percepita} & Adams & Ciò che do è in proporzione a ciò che ricevo? & Quali garanzie su carico, opportunità e riconoscimento concordare. \\
\bottomrule
\end{xltabular}}

Abbiamo scelto l'autodeterminazione al posto della gerarchia dei bisogni di Maslow, perché distingue \emph{perché} una persona si muove (e quindi come motivarla), non solo \emph{che cosa} le manca.

\subsection{Come è costruito}\label{sec:5.2}
\begin{itemize}
\item \textbf{12 domande}, una per dimensione in ciascuno di quattro momenti del cambiamento: annuncio, formazione, avvio operativo, uso dell'assistente AI. Sono scenari di lavoro, non domande astratte.
\item Le domande si ispirano a scale validate (Compeau e Higgins; Schwarzer e Jerusalem; Rigotti, Schyns e Mohr per l'autoefficacia; la Multidimensional Work Motivation Scale di Gagné e colleghi e gli orientamenti causali di Deci e Ryan per l'autodeterminazione; Adams per l'equità). Sono adattamenti: la loro validità va verificata nel contesto d'uso.
\item \textbf{Punteggio:} per ogni dimensione una media su 100 di quattro risposte di pari peso. Nell'autodeterminazione l'interesse vale 3 punti, il valore riconosciuto 2, la pressione 1, l'amotivazione 0. Nell'equità vale di più l'equilibrio, perché anche sentirsi avvantaggiati segnala criteri poco chiari.
\item \textbf{Soglia:} 50. Da 50 in su la dimensione è alta (1), sotto è bassa (0). Il codice a tre cifre segue l'ordine A-D-E. I punteggi tra 40 e 60 compaiono come \guillemotleft{}vicino alla soglia\guillemotright{} nella scheda personale.
\item Età (facoltativa) e termine di paragone dell'equità restano \textbf{fuori dal punteggio} e servono solo al colloquio o alle analisi aggregate.
\item Una demo funzionante del trivia è nel sito del progetto (link nella presentazione).
\end{itemize}

\subsection{Le otto personas e le tre categorie}\label{sec:5.3}
\textbf{Regola di classificazione (assunzione del gruppo):} è \textbf{Champion} chi ha autodeterminazione alta e almeno un'altra dimensione alta; è \textbf{Cinico} chi ha autodeterminazione ed equità basse; tutti gli altri sono \textbf{Fence sitters}. I nomi delle categorie riprendono le tre reazioni al cambiamento viste nel corso: campioni, seduti sulla staccionata, cinici.

{\footnotesize\renewcommand{\arraystretch}{1.25}
\begin{xltabular}{\linewidth}{@{}>{\raggedright\arraybackslash\hsize=0.631\hsize\hspace{0pt}}X>{\raggedright\arraybackslash\hsize=0.901\hsize\hspace{0pt}}X>{\raggedright\arraybackslash\hsize=0.793\hsize\hspace{0pt}}X>{\raggedright\arraybackslash\hsize=1.300\hsize\hspace{0pt}}X>{\raggedright\arraybackslash\hsize=1.209\hsize\hspace{0pt}}X>{\raggedright\arraybackslash\hsize=1.166\hsize\hspace{0pt}}X@{}}
\toprule
\textbf{Codice} & \textbf{Persona} & \textbf{Motivazione prevalente} & \textbf{Comportamento atteso} & \textbf{Che cosa proponiamo} & \textbf{Che cosa osserviamo per confermare} \\
\midrule
\endhead
\addlinespace[3pt]\multicolumn{6}{@{}l}{\textbf{Champions} (≈20\%, ≈170)} \\\addlinespace[1pt]
111 & Pioniere & intrinseca & Capace, convinto, trattato equamente: prende l'iniziativa e trascina i colleghi. & Un miglioramento scelto da lui; ruolo di utente chiave o formatore. & Esegue, mantiene i controlli, fa ripetere il compito a un collega. \\
110 & Talento in bilico & intrinseca & Ci crede, ma sente squilibrio su carichi o merito: può trattenere il contributo. & Accordi scritti su carico, responsabilità e riconoscimento. & Qualità e disponibilità ad aiutare dopo che gli accordi sono rispettati. \\
011 & Volenteroso in crescita & intrinseca & Vuole, ma gli manca la pratica. & Tutor e una seconda prova con meno aiuto; Champion in prospettiva. & Ripete il compito mantenendo i controlli. \\
\addlinespace[3pt]\multicolumn{6}{@{}l}{\textbf{Fence sitters} (≈70\%, ≈580)} \\\addlinespace[1pt]
101 & Esecutore in cerca di senso & estrinseca & Sa fare, ma non vede il perché: esegue e aspetta un beneficio concreto. & Una prova utile scelta insieme. & Prende iniziativa dopo aver visto il beneficio. \\
010 & Contributore da sostenere & intrinseca & Vuole, ma è insicuro e frenato da condizioni che sente inique. & Formazione accessibile, carico sostenibile, prova in coppia. & Partecipa e progredisce dopo gli accordi. \\
001 & Esploratore prudente & estrinseca & Teme l'errore; si muove solo se guidato. & Un compito piccolo, utile e guidato. & Avvia la prova; l'ostacolo si riduce. \\
\addlinespace[3pt]\multicolumn{6}{@{}l}{\textbf{Cinici} (≈10\%, ≈100)} \\\addlinespace[1pt]
100 & Esperto in cerca di garanzie & estrinseca & Esperto e critico: segnala difetti reali, può fare opposizione. & Carichi chiari, una prova che valorizzi l'esperienza; consultato per il feedback. & Partecipa dopo accordi rispettati; le critiche fondate cambiano il processo. \\
000 & Professionista in cerca di stabilità & amotivazione & Teme di perdere il ruolo; può dire sì e tornare ai vecchi flussi. & Impegni scritti con HR, una prova circoscritta con aiuto. & Risponde quando le condizioni sono mantenute. \\
\bottomrule
\end{xltabular}}

\textbf{Dove servono i Champions.} Circa 90 ruoli formali: 18 utenti chiave nel pilota ERP, 32 autori di automazioni, ≈40 referenti di funzione e sede, 2 posti nel team di governance (che possono coincidere con uno dei ruoli precedenti). Gli altri Champions sostengono i colleghi in modo informale (sessioni di supporto sull'AI, community, affiancamento al go-live delle ondate).

\textbf{Che cosa vale per tutti i Fence sitters:} workshop di aggiornamento, formazione nel tempo di lavoro, ascolto su richiesta. Non partecipano alle decisioni, ma sono sempre informati prima.

\textbf{I Cinici} vengono consultati per raccogliere feedback sulle soluzioni e trasferire ciò che sanno, senza ruolo decisionale: le loro critiche diventano informazione, e chi viene ascoltato resiste meno.

\subsection{I bias che stiamo assumendo}\label{sec:5.4}
La classificazione è un'assunzione. Per ogni persona dichiariamo le distorsioni che potrebbero falsarla: sia di chi risponde, sia nostre.

{\small\renewcommand{\arraystretch}{1.25}
\begin{xltabular}{\linewidth}{@{}>{\raggedright\arraybackslash\hsize=0.500\hsize\hspace{0pt}}X>{\raggedright\arraybackslash\hsize=0.350\hsize\hspace{0pt}}X>{\raggedright\arraybackslash\hsize=2.150\hsize\hspace{0pt}}X@{}}
\toprule
\textbf{Persona} & \textbf{Categoria} & \textbf{Bias che stiamo assumendo} \\
\midrule
\endhead
111 Pioniere & Champion & \textbf{Percezione selettiva} (nostra): vediamo l'entusiasmo e trascuriamo il rischio di uso non governato dell'AI. \textbf{Self-serving bias} (suo): si attribuisce i successi e sovrastima le proprie capacità. \\
110 Talento in bilico & Champion & \textbf{Ambivalenza}: il suo \guillemotleft{}sì, ma\guillemotright{} sull'equità rischia di essere letto come adesione piena. Assumiamo che lo squilibrio dipenda dal contesto, non dalla persona. \\
011 Volenteroso in crescita & Champion & \textbf{Effetto contrasto}: l'autoefficacia bassa può nascere dal confronto con colleghi più pratici. \textbf{Stereotipo} (nostro): associarlo ai più giovani. \\
101 Esecutore in cerca di senso & Fence sitter & \textbf{Dissonanza cognitiva}: si sente capace e trattato equamente, quindi può rifiutare la categoria e contestare il test. \textbf{Ambivalenza}: esegue, ma aspetta un beneficio. \\
010 Contributore da sostenere & Fence sitter & \textbf{Errore fondamentale di attribuzione} (nostro): leggere la sua attesa come disinteresse, quando può dipendere da formazione inaccessibile o carichi eccessivi. \textbf{Effetto contrasto} sull'equità. \\
001 Esploratore prudente & Fence sitter & \textbf{Percezione selettiva} (nostra): vediamo solo il timore dell'errore. \textbf{Stereotipo} (nostro): associarlo agli over 50 o alle funzioni con procedure consolidate. \\
100 Esperto in cerca di garanzie & Cinico & \textbf{Errore fondamentale di attribuzione} (nostro): la sua critica può segnalare difetti reali, non mancanza di volontà. \textbf{Stereotipo} (nostro): over 50, Finance, Operations, Procurement. La sua resistenza è spesso esplicita, quindi utile al confronto. \\
000 Professionista in cerca di stabilità & Cinico & \textbf{Self-serving bias} (suo): attribuisce il profilo al test e lo contesta. \textbf{Errore fondamentale di attribuzione} (nostro): il ritorno ai vecchi flussi può dipendere da un nuovo flusso poco praticabile. Rischio di resistenza implicita. \\
\bottomrule
\end{xltabular}}

Un bias vale per tutti: poiché il trivia è nominativo, alcune risposte saranno date come si pensa che l'azienda voglia. La restituzione (\hyperref[sec:7]{\S7}) e l'osservazione dei comportamenti durante il pilota servono anche a correggere questo effetto.

\subsection{Ipotesi sulla distribuzione}\label{sec:5.5}
Ci aspettiamo più Champions tra gli under 30 e in IT \& Digital Transformation, e più profili in cerca di garanzie e stabilità tra gli over 50 e nelle funzioni con procedure consolidate (Finance, Operations, Procurement). Sono \textbf{ipotesi da verificare, non etichette}: nessuna persona viene classificata per età o reparto, l'età è facoltativa e si legge solo in forma aggregata.

\subsection{Uso dei dati del trivia}\label{sec:5.6}
\begin{itemize}
\item \textbf{Obbligatorio e nominativo} (decisione del gruppo): serve ad assegnare i ruoli nel cambiamento (team pilota, utenti chiave, Champions delle automazioni) e a dimensionare formazione e supporto.
\item \textbf{Chi vede il profilo individuale:} solo la persona, HR e il responsabile del pilota.
\item \textbf{Mai per valutare} la prestazione o per decisioni sul personale (retribuzione, carriera, licenziamenti).
\item Le analisi si fanno su \textbf{dati aggregati}, con almeno 10 persone per gruppo.
\item Gli incarichi restano volontari: il profilo propone, la persona accetta (con l'accordo del responsabile).
\item \textbf{Conformità (da verificare con l'ufficio legale):} in Italia art. 8 dello Statuto dei lavoratori (divieto di indagini sulle opinioni non rilevanti per l'attitudine professionale) e art. 113 del Codice privacy; GDPR (artt. 5, 6, 13: finalità, minimizzazione, base giuridica, informativa); valutazione d'impatto (DPIA); in Germania il coinvolgimento del consiglio aziendale. Il CEO approva l'informativa prima della somministrazione.
\end{itemize}

\section{Analisi: problematiche e fonti della resistenza}\label{sec:6}

\subsection{I rischi per ciascuna tecnologia}\label{sec:6.1}

{\small\renewcommand{\arraystretch}{1.25}
\begin{xltabular}{\linewidth}{@{}>{\raggedright\arraybackslash\hsize=0.433\hsize\hspace{0pt}}X>{\raggedright\arraybackslash\hsize=1.866\hsize\hspace{0pt}}X>{\raggedright\arraybackslash\hsize=0.701\hsize\hspace{0pt}}X@{}}
\toprule
\textbf{Tecnologia} & \textbf{Che cosa può andare storto} & \textbf{Timori del Board toccati} \\
\midrule
\endhead
\textbf{ERP in cloud} & Vecchie procedure e file paralleli; dati incerti tra reparti; sapere degli esperti non scritto; carico e errori in aumento durante la migrazione. & Errori nella transizione · perdita di conoscenze \\
\textbf{GenAI} & Dati aziendali in strumenti personali (\emph{shadow AI}); risposte non verificate; paura che l'AI sostituisca il ruolo. & AI non governata · rischi cyber · resistenza \\
\textbf{Automazione} & Eccezioni note solo all'autore; errori del processo originario replicati in serie; ore liberate che diventano pressione invece che sviluppo. & Perdita di conoscenze · errori · resistenza \\
\bottomrule
\end{xltabular}}

\subsection{Le problematiche trasversali emerse dai casi}\label{sec:6.2}
\begin{enumerate}
\item \textbf{Paura di perdere il ruolo} o l'autonomia (HolidayCo, Qantas).
\item \textbf{Assimilazione incompleta} delle tecnologie, con strumenti paralleli (Robey et al., Lidl).
\item \textbf{Perdita della conoscenza di processo} quando le attività passano ai bot (studi sull'automazione robotica dei processi).
\item \textbf{Uso non governato di AI e agenti} (rapporti Okta e Gartner sulla \emph{shadow AI}).
\item \textbf{Formazione senza un nuovo ruolo}: le competenze non si traducono in responsabilità e il tempo liberato non viene reinvestito (TCS, Accenture, IBM).
\item \textbf{Scarso coordinamento tra reparti e poco ascolto} dei dipendenti (HolidayCo, SACE).
\item \textbf{Sovraccarico durante la transizione}, soprattutto con l'ERP.
\end{enumerate}

Le tre tecnologie hanno quindi problemi specifici, ma condividono cinque dimensioni: competenze, comunicazione, collaborazione tra funzioni, riconoscimento dei contributi, sostenibilità dei carichi. Il cambiamento non si può gestire come un progetto solo tecnologico.

\subsection{Da dove viene la resistenza}\label{sec:6.3}
\begin{itemize}
\item \textbf{Individuale:} abitudini, bisogno di sicurezza, paura per il posto e dell'ignoto, percezione selettiva.
\item \textbf{Organizzativa:} inerzia della struttura e del gruppo, minaccia alle competenze, al potere e alle risorse.
\item \textbf{Come si manifesta:} in modo \textbf{esplicito} (opposizione dichiarata) oppure \textbf{implicito e differito} (mancata collaborazione, scarso uso degli strumenti nuovi, procedure parallele). L'assenza di opposizione non equivale ad adesione.
\end{itemize}

\subsection{Che cosa ne ricaviamo per il piano}\label{sec:6.4}

{\small\renewcommand{\arraystretch}{1.25}
\begin{xltabular}{\linewidth}{@{}>{\raggedright\arraybackslash\hsize=0.357\hsize\hspace{0pt}}X>{\raggedright\arraybackslash\hsize=1.643\hsize\hspace{0pt}}X@{}}
\toprule
\textbf{Problema} & \textbf{Risposta nel piano} \\
\midrule
\endhead
Paura per il ruolo & Nessun taglio legato alle automazioni, ore reinvestite, ruoli nuovi per chi si forma (\hyperref[sec:11]{\S11}, \hyperref[sec:13.6]{\S13.6}) \\
Assimilazione incompleta & Pilota con utenti chiave, sandbox, competenza come soglia di go-live, KPI sugli strumenti paralleli (\hyperref[sec:9]{\S9}, \hyperref[sec:12]{\S12}) \\
Perdita di conoscenza & ERP core team, manuale della linea, sostituto per ogni automazione, mentoring senior-junior (\hyperref[sec:8]{\S8}, \hyperref[sec:11]{\S11}, \hyperref[sec:12]{\S12}) \\
Shadow AI & Strumento aziendale per tutti, regole chiare, registro che dà riconoscimento (\hyperref[sec:10]{\S10}, \hyperref[sec:11]{\S11}, \hyperref[sec:13]{\S13}) \\
Formazione senza ruolo & Ogni percorso porta a un incarico (utente chiave, autore di automazioni, referente) (\hyperref[sec:12]{\S12}) \\
Poco coordinamento e ascolto & Team di prodotto come unità del cambiamento, referenti, Cinici consultati, ciclo ``ci avete detto / abbiamo cambiato'' (\hyperref[sec:7]{\S7}, \hyperref[sec:8]{\S8}) \\
Sovraccarico & Tempo protetto, sostituti temporanei, pulse mensile sul carico (\hyperref[sec:12]{\S12}, \hyperref[sec:15]{\S15}, \hyperref[sec:16]{\S16}) \\
\bottomrule
\end{xltabular}}

\section{Feedback: restituire il profilo e leggere le reazioni}\label{sec:7}

\subsection{La restituzione in quattro passaggi}\label{sec:7.1}

{\small\renewcommand{\arraystretch}{1.25}
\begin{xltabular}{\linewidth}{@{}>{\raggedright\arraybackslash\hsize=0.265\hsize\hspace{0pt}}X>{\raggedright\arraybackslash\hsize=1.735\hsize\hspace{0pt}}X@{}}
\toprule
\textbf{Quando} & \textbf{Che cosa succede} \\
\midrule
\endhead
\textbf{Prima della plenaria} (M2) & Ognuno riceve la \textbf{scheda della propria persona}: profilo, ruolo nel cambiamento, supporto previsto, che cosa ci si aspetta, eventuale segnalazione \guillemotleft{}vicino alla soglia\guillemotright{}. \\
\textbf{In plenaria} & Si presentano i risultati aggregati e si pone una domanda: \guillemotleft{}Ti riconosci in questo profilo?\guillemotright{} \\
\textbf{Su richiesta} & Colloquio con HR: il profilo si può correggere. Non sono previsti colloqui pianificati per tutti. \\
\textbf{A M10} & Trivia ripetuto e comportamenti osservati aggiornano i profili e i ruoli. \\
\bottomrule
\end{xltabular}}

Il profilo è un'ipotesi di lavoro, non un giudizio. La prima classificazione serve subito (per scegliere il team pilota), la restituzione la conferma o la corregge, e la verifica si ripete durante il pilota: la classificazione è \textbf{ricorsiva}.

\subsection{Le reazioni che leggiamo}\label{sec:7.2}

{\small\renewcommand{\arraystretch}{1.25}
\begin{xltabular}{\linewidth}{@{}>{\raggedright\arraybackslash\hsize=0.456\hsize\hspace{0pt}}X>{\raggedright\arraybackslash\hsize=1.221\hsize\hspace{0pt}}X>{\raggedright\arraybackslash\hsize=1.324\hsize\hspace{0pt}}X@{}}
\toprule
\textbf{Reazione} & \textbf{Come si riconosce} & \textbf{Che cosa facciamo} \\
\midrule
\endhead
\textbf{Resistenza esplicita} & Contesta apertamente il profilo. & È utile: apre il confronto e rivela che cosa non convince. \\
\textbf{Resistenza implicita} & Dice sì, poi non partecipa alle prove o torna ai vecchi flussi. & La vediamo nei dati (strumenti paralleli, adozione dell'ERP, uso dell'AI). \\
\textbf{Dissonanza cognitiva} & Si sente innovatore ma risulta Fence sitter, e scredita il test o dice che il suo reparto è diverso. & Spieghiamo che il profilo descrive condizioni, non valore; offriamo una prova concreta. \\
\textbf{Bias di chi risponde} & Si attribuisce i successi e dà la colpa al test; risponde come vuole l'azienda. & Confronto con i comportamenti osservati. \\
\textbf{Bias nostri} & Leggere un'attesa come mancanza di volontà; legare il profilo a età o reparto. & Prima di attribuire una causa alla persona, si controlla il contesto (formazione, carico, strumenti). \\
\textbf{Ambivalenza} & \guillemotleft{}Sì, ma\guillemotright{}: adesione con condizioni. & Le condizioni si ascoltano e diventano accordi verificabili. \\
\bottomrule
\end{xltabular}}

\subsection{Dal feedback all'azione}\label{sec:7.3}
Ogni obiezione concreta diventa una condizione da verificare con il responsabile. Il Change \& Risk Office pubblica ogni due mesi un resoconto \textbf{``ci avete detto / abbiamo cambiato''}, che dice anche quali osservazioni non sono state accolte e perché. Una comunicazione coerente nel tempo è essenziale: promesse smentite colpiscono soprattutto chi ha già un'equità percepita bassa.

\section{Azione · Governance e struttura del cambiamento}\label{sec:8}

\subsection{Principio}\label{sec:8.1}
La struttura hard resta funzionale. I team di prodotto esistono già e sono gli unici stabili. \textbf{Tutto ciò che aggiungiamo è temporaneo} e si scioglie a fine programma, passando le attività a IT, alle linee e ai profili assunti. Separiamo con chiarezza chi \textbf{decide} (governance), chi \textbf{coordina} (Change \& Risk Office) e chi \textbf{esegue} (team di prodotto).

\subsection{Le strutture}\label{sec:8.2}

{\small\renewcommand{\arraystretch}{1.25}
\begin{xltabular}{\linewidth}{@{}>{\raggedright\arraybackslash\hsize=0.521\hsize\hspace{0pt}}X>{\raggedright\arraybackslash\hsize=1.306\hsize\hspace{0pt}}X>{\raggedright\arraybackslash\hsize=1.734\hsize\hspace{0pt}}X>{\raggedright\arraybackslash\hsize=0.439\hsize\hspace{0pt}}X@{}}
\toprule
\textbf{Struttura} & \textbf{Composizione} & \textbf{Che cosa fa} & \textbf{Durata} \\
\midrule
\endhead
\textbf{Team di governance} · decide & CEO (sponsor), CFO, COO, CIO, direttore HR, CISO, 2 Champions (uno del pilota, uno delle automazioni), consulenti KPMG in ruolo consultivo senza voto. ≈10 persone. & Fissa le soglie e decide i go/no-go, approva la classificazione del rischio delle automazioni, il piano di formazione e il budget. Riunione mensile. Decide, non esegue. & M1-M22, poi guida la chiusura \\
\textbf{Change \& Risk Office} · coordina & 4-6 persone, in staff al CEO. & Raccordo unico tra pilota, automazioni, formazione e recognition; ogni mese controlla con l'IT spesa e uso dell'AI (console aggregata); report alla governance; ciclo ``ci avete detto / abbiamo cambiato''. Nessuna autorità gerarchica sulle linee. & M1-M22 \\
\textbf{Referenti di funzione e sede} & ≈40 Champions, 10\% del tempo formalizzato e un obiettivo nella valutazione. & Portano la voce dei colleghi, animano la community, raccolgono problemi e proposte. & M1-M22 \\
\textbf{Team di prodotto del pilota} · esegue & Il team già esistente della linea scelta (≈60 persone NexaCore) più 5 consulenti KPMG per la sola durata del pilota. & Porta la linea sul nuovo ERP. Il responsabile di linea risponde del servizio ai clienti e dei risultati. & Stabile; KPMG solo M2-M10 \\
\textbf{ERP core team} & 6-8 persone: 3-4 utenti chiave del pilota, 2 dall'IT, 1 referente dati e autorizzazioni. Riporta al CIO (tecnica) e ai process owner (processo). & Configurazioni, rilasci, profili di accesso, manuale interno; forma le linee delle ondate. Prende il posto di KPMG: senza, la conoscenza uscirebbe con i consulenti. & M9-M22 \\
\textbf{Process owner} & Uno per processo end-to-end (ordine-incasso, acquisto-pagamento), nelle linee. & Rispondono del risultato operativo, non del progetto. & Stabile \\
\bottomrule
\end{xltabular}}

\textbf{Regola anti-ambiguità.} Le persone con un ruolo nel progetto hanno due riferimenti (reparto e progetto): è il rischio noto delle strutture a matrice. Per evitarlo, \textbf{un solo capo valuta la persona}, e ruolo nel progetto e percentuale di tempo dedicata sono messi per iscritto.

\textbf{Management partecipativo e partecipazione rappresentativa.} I Champions siedono in governance, i referenti coprono tutte le funzioni e sedi, i Cinici sono consultati: chi è toccato dal cambiamento ha voce nelle decisioni, anche se non le prende da solo.

\subsection{Il ruolo di KPMG}\label{sec:8.3}
\begin{itemize}
\item \textbf{Perché in governance.} Soglie di go/no-go e classificazione del rischio sono decisioni tecniche che si prendono una volta: chi ha visto altre transizioni ERP sa dove si rompono. Il parere su cyber e AI serve \textbf{prima} della decisione.
\item \textbf{Chi.} Un partner o senior manager al comitato mensile; l'esperto cyber al Change \& Risk Office quando si trattano accessi, classificazione del rischio e casi d'uso AI. Non sono le stesse persone che configurano il sistema.
\item \textbf{Limiti.} Ruolo \textbf{consultivo, senza voto}: le decisioni restano a NexaCore. Le soglie si fissano a M1, prima del go-live, non mentre lo si valuta.
\item \textbf{Conflitto di interessi.} KPMG configura il sistema e consiglia sul go/no-go dello stesso sistema. Mitigazione: persone diverse, soglie scritte in anticipo, verbale che riporta \textbf{parere di KPMG e decisione del team} come voci separate.
\item \textbf{Presenza nel tempo.} Nel team del pilota: manager al 50\%, 3 consulenti ERP, esperto cyber part-time; presenza piena M2-M7, in calo da M8, uscita a M10. In governance: presenza piena fino a M10, poi una giornata al mese fino a M22, più una revisione a fine ondata. Nelle ondate: adattamento ai paesi (fisco, lingue). Nell'ERP core team: nessun consulente.
\end{itemize}

\subsection{Quale ruolo a quale Champion}\label{sec:8.4}
Per assegnare i Champions ai ruoli usiamo come \textbf{criterio di affinità} (non come regola di selezione) i bisogni di McClelland, chiesti al colloquio di assegnazione con HR e non misurati nel trivia:
\begin{itemize}
\item \textbf{Potere} (influenzare gli altri) → team di governance: portavoce della propria funzione davanti a CEO e KPMG.
\item \textbf{Affiliazione} (relazioni con i colleghi) → utente chiave ERP: forma, affianca, raccoglie i problemi.
\item \textbf{Realizzazione} (obiettivi sfidanti con riscontro rapido) → autore di automazioni: risultato misurabile (ore liberate, errori evitati).
\end{itemize}

\section{Azione · Pilota ERP su una linea di prodotto}\label{sec:9}

\subsection{Perché una linea di prodotto}\label{sec:9.1}
Il nuovo ERP non si accende in tutta l'azienda nello stesso giorno. Si parte da \textbf{una sola linea di prodotto}, perché una linea attraversa tutti i reparti, dalla vendita all'acquisto, dalla produzione all'amministrazione all'assistenza: il pilota collauda l'intero percorso dall'ordine all'incasso, cioè proprio i passaggi tra reparti dove nascono gli errori della transizione.
\begin{itemize}
\item Un pilota su \textbf{una sola funzione} collauderebbe un pezzo isolato, senza i passaggi: poco rappresentativo.
\item Un pilota su \textbf{un'intera sede} coinvolgerebbe più linee e molte più persone insieme: più rischio per i clienti.
\end{itemize}

\subsection{Quale linea}\label{sec:9.2}
Si sceglie il \textbf{team di prodotto già esistente con il maggior numero di Champions} (dati del trivia). Chi è favorevole al cambiamento assorbe meglio i disagi, quindi la linea continua a servire i clienti. Assumiamo, in modo ottimistico, che in quel team i Champions siano distribuiti tra le funzioni, che la linea non abbia picchi di domanda durante il pilota e che la maggior parte delle persone lavori in una sede e in un paese (una normativa, una lingua, meno integrazioni).

\textbf{Conseguenza dichiarata:} i risultati del pilota sono un \textbf{benchmark ottimistico}, non una promessa. Le ondate hanno soglie più prudenti.

\textbf{Procedura:} i dipendenti compilano il trivia (M1) → classificazione → HR e IT contano i Champions per team → la governance approva il primo in classifica (scelta provvisoria, confermata dopo la restituzione di M2) → il responsabile di linea con KPMG definisce la composizione → il CEO apre il pilota.

\subsection{Chi partecipa}\label{sec:9.3}

{\small\renewcommand{\arraystretch}{1.25}
\begin{xltabular}{\linewidth}{@{}>{\raggedright\arraybackslash\hsize=0.469\hsize\hspace{0pt}}X>{\raggedright\arraybackslash\hsize=0.389\hsize\hspace{0pt}}X>{\raggedright\arraybackslash\hsize=0.228\hsize\hspace{0pt}}X>{\raggedright\arraybackslash\hsize=2.914\hsize\hspace{0pt}}X@{}}
\toprule
\textbf{Team pilota (≈65 persone)} & \textbf{Persone} & \textbf{\%} & \textbf{Ruolo} \\
\midrule
\endhead
Champions & 18 & 28\% & Utenti chiave, almeno uno per funzione: provano i processi, formano i colleghi, raccolgono i problemi \\
Fence sitters & 36 & 55\% & Ricevono istruzioni e formazione già pronte; non decidono, partecipano ai workshop di aggiornamento \\
Cinici & 6 & 9\% & 2-3 tenuti vicini nei test e nelle revisioni degli errori, consultati senza ruolo decisionale \\
KPMG & 5 & 8\% & Manager al 50\%, 3 consulenti ERP, esperto cyber part-time: configurano, presidiano il rischio cyber, trasferiscono competenze \\
\bottomrule
\end{xltabular}}

\subsection{Come si svolge (M2-M10)}\label{sec:9.4}

{\small\renewcommand{\arraystretch}{1.25}
\begin{xltabular}{\linewidth}{@{}>{\raggedright\arraybackslash\hsize=0.207\hsize\hspace{0pt}}X>{\raggedright\arraybackslash\hsize=1.793\hsize\hspace{0pt}}X@{}}
\toprule
\textbf{Periodo} & \textbf{Che cosa succede} \\
\midrule
\endhead
\textbf{M2-M3} & Si misura la linea oggi (errori, tempi, consegne puntuali). Workshop sui processi standard dell'ERP con KPMG e utenti chiave: che cosa adottare così com'è e che cosa adattare. \\
\textbf{M3-M5} & KPMG configura a cicli brevi; gli utenti chiave verificano ogni ciclo. Pulizia e migrazione dei dati; profili di accesso (chi può fare cosa). \\
\textbf{M6} & \textbf{Prova generale} su casi reali; gli utenti chiave formano i colleghi. \textbf{Go/no-go del go-live}: si accende solo se test, dati e competenza degli utenti chiave superano le soglie fissate a M1. \\
\textbf{M7-M8} & \textbf{Go-live}. Il vecchio sistema resta disponibile per \textbf{4 settimane} come rete di sicurezza. Affiancamento quotidiano; revisione settimanale degli errori senza ricerca di colpevoli. Train-the-trainer degli utenti chiave (M7-M10). \\
\textbf{M8} & KPMG comincia a ridurre la presenza: gli utenti chiave conducono, i consulenti controllano. \\
\textbf{M9} & Si spegne il vecchio sistema. Nasce l'\textbf{ERP core team}. \\
\textbf{M10} & Misura finale e T1 (con trivia ripetuto). \textbf{Go/no-go dell'estensione}: si passa alle altre linee, e si sblocca la tranche 2, solo se continuità, errori e commitment sono entro le soglie. \textbf{KPMG esce dal team pilota.} \\
\bottomrule
\end{xltabular}}

Il pilota dura circa nove mesi, coerente con i 6-10 mesi tipici del primo avvio di un ERP in cloud.

\subsection{Trasferimento delle competenze}\label{sec:9.5}
Il compito di KPMG è rendere NexaCore autonoma. Ogni utente chiave lavora in coppia con un consulente fin dalla configurazione; dall'avvio i ruoli si invertono. Processi e soluzioni ai problemi finiscono nel \textbf{manuale interno della linea}. Robey et al. distinguono due barriere: la \textbf{configurazione} (la copre KPMG) e l'\textbf{assimilazione} nel lavoro quotidiano (la presidiano gli utenti chiave).

\subsection{Le ondate (M11-M22)}\label{sec:9.6}
\begin{itemize}
\item 4 linee, ≈4 mesi ciascuna, una ogni ≈2,5 mesi: avvio a \textbf{M11, M13, M16, M18}; l'ultima chiude a M22.
\item Al massimo \textbf{due linee in transizione insieme}, mai tre.
\item Le formano gli \textbf{utenti chiave del pilota} e l'ERP core team; KPMG adatta a fisco e lingue dei paesi.
\item Ogni ondata ha il suo go/no-go e un premio una tantum alla linea se le soglie di continuità tengono.
\item A ogni ondata si integrano nell'ERP le automazioni migliori della linea (\hyperref[sec:11]{\S11}).
\end{itemize}

\section{Azione · Intelligenza artificiale generativa per tutti}\label{sec:10}

\subsection{Perché a tutti e subito}\label{sec:10.1}
Il pilota ERP tocca ≈65 persone; le altre 785 non possono aspettare due anni. Inoltre, se lo strumento aziendale manca, le persone continuano a usare quelli personali: la \emph{shadow AI} si sostituisce, non si vieta soltanto.

\subsection{Il contratto enterprise: i requisiti}\label{sec:10.2}
Si stipula un contratto enterprise con un fornitore di modelli linguistici. Non scegliamo un'azienda precisa: indichiamo i requisiti che il contratto deve garantire.
\begin{itemize}
\item I dati inseriti e generati \textbf{restano di NexaCore}: non vengono ceduti a terzi né usati per addestrare i modelli del fornitore.
\item Account aziendali nominativi, amministrati centralmente dall'IT, con gestione degli accessi.
\item \textbf{Console di amministrazione} con statistiche d'uso (utenti attivi, volumi, funzionalità, consumi) lette \textbf{solo in forma aggregata} (per funzione, sede, profilo), \textbf{mai sul singolo dipendente}.
\item Prezzi a volume e tetti di spesa contrattuali.
\end{itemize}
I requisiti vanno verificati nelle condizioni effettive del contratto.

\subsection{Tre profili d'accesso}\label{sec:10.3}

{\small\renewcommand{\arraystretch}{1.25}
\begin{xltabular}{\linewidth}{@{}>{\raggedright\arraybackslash\hsize=0.404\hsize\hspace{0pt}}X>{\raggedright\arraybackslash\hsize=1.168\hsize\hspace{0pt}}X>{\raggedright\arraybackslash\hsize=1.428\hsize\hspace{0pt}}X@{}}
\toprule
\textbf{Profilo} & \textbf{Uso previsto} & \textbf{Accesso e limiti} \\
\midrule
\endhead
\textbf{Base} & Scrittura e revisione di testi, sintesi, attività semplici. & Modelli per le attività ordinarie; disponibilità contenuta per le funzioni più costose. \\
\textbf{Intermedio} & Analisi di documenti e dati, report, attività in più passaggi. & Accesso più ampio, con budget mensile limitato per modelli potenti e analisi complesse. \\
\textbf{Avanzato} & Sviluppo e verifica di automazioni, supporto alla programmazione e ai progetti ERP. & Strumenti specialistici e budget maggiore, legato ai progetti autorizzati. \\
\bottomrule
\end{xltabular}}

Il profilo dipende da \textbf{competenze dimostrate e necessità del ruolo}, non dalla categoria motivazionale: un Champion non riceve automaticamente il profilo avanzato. I permessi sull'ERP e sui dati restano separati e seguono la classificazione del rischio. Scenario di dimensionamento: 600 base, 200 intermedi, 50 avanzati (da confermare).

\subsection{Tempi}\label{sec:10.4}

{\small\renewcommand{\arraystretch}{1.25}
\begin{xltabular}{\linewidth}{@{}>{\raggedright\arraybackslash\hsize=0.288\hsize\hspace{0pt}}X>{\raggedright\arraybackslash\hsize=0.496\hsize\hspace{0pt}}X>{\raggedright\arraybackslash\hsize=2.216\hsize\hspace{0pt}}X@{}}
\toprule
\textbf{Periodo} & \textbf{Fase} & \textbf{Attività} \\
\midrule
\endhead
M1 & Contratto & Scelta del fornitore e firma del contratto (prime 4 settimane); regole su dati, profili, budget, responsabilità. \\
M2-M3 & Referenti & Configurazione e collaudo dell'ambiente; account e formazione avanzata dei 32 Champions delle automazioni. \\
M4-M6 & Diffusione & Account e formazione progressiva in tutte le 8 funzioni; demo, sessioni di supporto, libreria di esempi validati. \\
M6-M22 & Consolidamento & Monitoraggio mensile di utilizzo, qualità e costi; aggiornamento dei profili; supporto alle ondate ERP. \\
\bottomrule
\end{xltabular}}

\subsection{Controllo dei costi}\label{sec:10.5}
I costi dell'AI crescono con l'uso, soprattutto per le automazioni (si paga a token, cioè per quantità di testo inviato e generato). Controlli previsti: tetti mensili per profilo con avviso all'80\% e blocco al 100\%; crediti per progetto per gli avanzati; il modello più economico come scelta predefinita; costo stimato di ogni automazione nel registro; controllo mensile del Change \& Risk Office.

\section{Azione · Automazioni dei Champions}\label{sec:11}

\subsection{Che cosa e chi}\label{sec:11.1}
Mentre il pilota porta l'ERP su una linea, nel resto dell'azienda parte un secondo cantiere, più leggero: \textbf{i Champions automatizzano le attività ripetitive del proprio lavoro} (copiare dati tra file, confrontare elenchi, preparare lo stesso report ogni settimana, smistare richieste).
\begin{itemize}
\item \textbf{≈32 Champions}, 4 per ciascuna delle 8 funzioni, \textbf{fuori dal team pilota}, perché chi implementa l'ERP non abbia un secondo compito e il cambiamento arrivi ovunque.
\item L'incarico è \textbf{volontario}: lo riceve chi è Champion, ha completato la formazione e lo accetta, con l'accordo del responsabile.
\item \textbf{Mezza giornata a settimana} di tempo protetto (≈10\%); account con profilo avanzato.
\end{itemize}

\subsection{Le regole}\label{sec:11.2}
\begin{itemize}
\item \textbf{Solo strumenti aziendali.}
\item \textbf{Classificazione del rischio}, decisa dalla governance (KPMG dà parere senza voto): stabilisce che cosa si può automatizzare e con quali approvazioni. Dati personali e pagamenti sono esclusi dal perimetro iniziale.
\item \textbf{Registro} di ogni automazione: processo, dati usati, autore, sostituto, ore risparmiate, costo stimato.
\item \textbf{Sostituto obbligatorio} e istruzioni scritte: la conoscenza non resta in una sola testa.
\item \textbf{Controllo umano} dell'esito finché l'automazione non ha dimostrato di funzionare.
\item \textbf{Ore reinvestite}, nessun taglio di posti: dirlo apertamente è la condizione perché i colleghi accettino le automazioni.
\item Un'automazione registrata con istruzioni e sostituto genera un \textbf{riconoscimento automatico} sulla piattaforma di recognition: registrare conviene, e l'AI usata di nascosto viene alla luce.
\end{itemize}

\subsection{Il collegamento con l'ERP}\label{sec:11.3}
Ogni mese, in ogni funzione, Champions delle automazioni e utenti chiave dell'ERP si incontrano per un'ora (conduzione a rotazione, sintesi al responsabile di funzione). Ordine del giorno fisso: automazioni nuove e livello di rischio; che cosa cambia nei processi con l'ERP; automazioni da integrare o fermare; problemi aperti. Così nessuno automatizza ciò che l'ERP cambierà, e le automazioni migliori diventano candidate all'integrazione nel nuovo sistema.

\subsection{Tempi}\label{sec:11.4}

{\small\renewcommand{\arraystretch}{1.25}
\begin{xltabular}{\linewidth}{@{}>{\raggedright\arraybackslash\hsize=0.213\hsize\hspace{0pt}}X>{\raggedright\arraybackslash\hsize=1.787\hsize\hspace{0pt}}X@{}}
\toprule
\textbf{Mese} & \textbf{Tappa} \\
\midrule
\endhead
M2-M3 & Formazione: 24 ore di AI avanzata, poi laboratori con specialisti esterni \\
M3 & La governance fissa la \textbf{classificazione del rischio} e il registro; account attivi \\
M4 & \textbf{Prime automazioni} e prime call mensili \\
M6 & \textbf{Revisione} delle prime automazioni; la governance decide se allargare il perimetro \\
M10 & Le automazioni efficaci entrano nell'ERP della linea pilota \\
M11-M22 & A ogni ondata si integrano nell'ERP le automazioni della linea \\
\bottomrule
\end{xltabular}}

Uno dei Champions delle automazioni siede nel team di governance e porta l'esperienza delle automazioni nelle decisioni sulla formazione per tutti.

\section{Azione · Formazione e conservazione della conoscenza}\label{sec:12}

\subsection{Percorsi}\label{sec:12.1}

{\small\renewcommand{\arraystretch}{1.25}
\begin{xltabular}{\linewidth}{@{}>{\raggedright\arraybackslash\hsize=0.656\hsize\hspace{0pt}}X>{\raggedright\arraybackslash\hsize=2.274\hsize\hspace{0pt}}X>{\raggedright\arraybackslash\hsize=0.515\hsize\hspace{0pt}}X>{\raggedright\arraybackslash\hsize=0.556\hsize\hspace{0pt}}X@{}}
\toprule
\textbf{Destinatari} & \textbf{Contenuto} & \textbf{Durata} & \textbf{Quando} \\
\midrule
\endhead
Utenti ERP della linea & Nuove procedure, passaggi tra reparti; sandbox e prova in coppia con un utente chiave & secondo il ruolo & prima di ogni go-live \\
Utenti chiave ERP (18 nel pilota) & Configurazione, integrazioni, problemi tecnici; poi \textbf{train-the-trainer} certificato & in coppia con KPMG M2-M10 & M2-M10 \\
AI · livello base & Funzionamento, compiti semplici, dati utilizzabili, verifica delle risposte & 4 ore & M4-M6 \\
AI · livello intermedio & Uso quotidiano, analisi di documenti e dati, qualità dei risultati & 8 ore & M4-M6 \\
AI · livello avanzato & Strumenti specialistici, agenti, supporto ai colleghi & 24 ore & M2-M3 (Champions) \\
Champions delle automazioni & Dopo l'avanzato: specialisti esterni della piattaforma → affiancamento sulle prime automazioni con IT → trasferimento, quando diventano referenti & laboratori & M2-M6 \\
Tutti gli 850 & \textbf{Cyber essentials} con simulazioni di phishing e richiami periodici & breve & M3-M5 \\
\bottomrule
\end{xltabular}}

Il tempo interno di formazione sull'AI è ≈5.200 ore (600 × 4 + 200 × 8 + 50 × 24). La formazione si svolge \textbf{nell'orario di lavoro} (tempo protetto).

\subsection{Presidi della conoscenza}\label{sec:12.2}
\begin{itemize}
\item \textbf{Sandbox ERP} permanente con dati fittizi: ci si esercita e si sbaglia senza toccare i dati reali.
\item \textbf{Competenza come soglia di go-live}: la prova generale di M6 non si supera se gli utenti chiave non sono autonomi.
\item \textbf{Knowledge base} (procedure, microvideo, manuale della linea) alimentata dalla community, più una libreria di esempi AI validati.
\item \textbf{Mentoring reciproco senior-junior}: i 180 over 50 trasferiscono la conoscenza tacita di processo, i più giovani le competenze digitali.
\item \textbf{ERP core team} (\hyperref[sec:8.2]{\S8.2}) e \textbf{sostituto per ogni automazione} (\hyperref[sec:11.2]{\S11.2}).
\end{itemize}

\subsection{Perché funziona: le leve}\label{sec:12.3}
\begin{itemize}
\item \textbf{Autoefficacia (Bandura):} si costruisce con la padronanza (sandbox, prova in coppia) e con l'esempio di colleghi simili (utenti chiave che formano).
\item \textbf{Equità (Adams):} la formazione si fa in orario e chi si carica di più riceve un'indennità (\hyperref[sec:13.6]{\S13.6}); altrimenti chi si forma percepirebbe uno scambio iniquo.
\item \textbf{Obiettivi (Locke):} le soglie del go-live sono specifiche e note da M1, con riscontri settimanali.
\item \textbf{Aspettativa (Vroom):} chi si forma vede che lo sforzo porta al risultato e il risultato a un riconoscimento o a un ruolo.
\item \textbf{Vittoria rapida (Kotter):} il team con più Champions mostra presto che il cambiamento funziona.
\end{itemize}

\section{Azione · Cultura, comunicazione e recognition}\label{sec:13}

\subsection{Comunicare partendo dal perché}\label{sec:13.1}
Il CEO apre il programma con il \textbf{cerchio d'oro} di Sinek: prima \textbf{perché} si cambia (proteggere clienti, dati e persone mentre l'azienda cresce), poi \textbf{come} (pilota, formazione, ascolto), solo alla fine \textbf{che cosa} (ERP, AI, automazioni).

\subsection{Cinque comportamenti osservabili}\label{sec:13.2}
I valori (innovazione, attenzione al dettaglio, orientamento al risultato) non si osservano e non si premiano; i comportamenti sì. Scelti in workshop a M1-M2 con Champions e responsabili, presentati a tutti, rivisti una volta l'anno.

{\small\renewcommand{\arraystretch}{1.25}
\begin{xltabular}{\linewidth}{@{}>{\raggedright\arraybackslash\hsize=0.639\hsize\hspace{0pt}}X>{\raggedright\arraybackslash\hsize=1.824\hsize\hspace{0pt}}X>{\raggedright\arraybackslash\hsize=0.537\hsize\hspace{0pt}}X@{}}
\toprule
\textbf{Comportamento} & \textbf{Che cosa si vede} & \textbf{Timore presidiato} \\
\midrule
\endhead
\textbf{1. Decido con i dati} & Uso il dato dell'ERP invece del file personale; chiudo l'attività a sistema il giorno stesso & Errori nella transizione \\
\textbf{2. Uso l'AI entro le regole} & Strumento aziendale, profilo assegnato, verifica dell'output & AI non governata · cyber \\
\textbf{3. Documento ciò che so} & Procedure, microvideo, automazione registrata con autore e sostituto & Perdita di conoscenze \\
\textbf{4. Segnalo errori e incidenti} & Segnalazione in giornata, revisione settimanale senza colpevoli & Errori · cyber \\
\textbf{5. Aiuto un collega} & Affiancamento al go-live, risposte nella community, formazione a un collega & Resistenza \\
\bottomrule
\end{xltabular}}

Cinque e non dieci: con troppe opzioni l'etichetta si sceglie a caso e la misura perde significato. I comportamenti entrano anche nelle checklist dell'ERP, nell'onboarding, nei colloqui di sviluppo e nel registro delle automazioni.

\subsection{La piattaforma di recognition}\label{sec:13.3}
\begin{itemize}
\item \textbf{Messaggio pubblico tra pari} con il \textbf{tag} obbligatorio di uno dei 5 comportamenti: da collega a collega, da responsabile a team, tra funzioni e paesi. Esempio: \emph{[Documento ciò che so] \guillemotleft{}Giulia ha scritto la procedura di chiusura degli ordini e l'ha messa nella knowledge base: ci ha risparmiato la telefonata di ogni lunedì.\guillemotright{}}
\item Il tag dice a chi riceve \textbf{che cosa} ha fatto di utile, e misura quali comportamenti si diffondono per funzione e paese senza un questionario.
\item \textbf{Avvio in due tempi:} M3-M5 solo messaggi e tag; da \textbf{M6} punti e catalogo premi (prima nel pilota e tra i Champions delle automazioni, poi con le ondate).
\item \textbf{Punti:} ogni persona ha un budget mensile da donare (es. 100, al massimo 30 allo stesso collega): il potere di riconoscere è distribuito, non concentrato nei capi. Bonus di traguardo al team a fine ondata.
\item \textbf{Nessun punto per il volume d'uso} degli strumenti, che premierebbe l'uso dell'AI fuori regole. Si premiano comportamento e contributo alla conoscenza.
\item \textbf{Catalogo a tre fasce} (buoni e gadget, accessori, elettronica), uguale nei 5 paesi con valore adeguato al costo della vita locale; punti con scadenza a 18-24 mesi, fuori dalla retribuzione fissa.
\item \textbf{Equità e anti-abuso:} niente auto-assegnazione, tetti per coppia, avvisi sugli scambi reciproci, distribuzione pubblicata ogni trimestre.
\item \textbf{Copertura:} una vista segnala chi non dà né riceve riconoscimenti (spesso Fence sitters e lavoratori da remoto), per un contatto mirato, non per una classifica.
\end{itemize}

\subsection{Perché è progettata così}\label{sec:13.4}
\begin{itemize}
\item \textbf{Rinforzo positivo (Skinner):} un comportamento seguito da una conseguenza positiva tende a ripetersi, se il riconoscimento è legato a un comportamento preciso (il tag) e arriva subito (il registro genera il riconoscimento in automatico). M3-M5: rinforzo continuo per far imparare; da M6 punti non prevedibili (rinforzo intermittente) per mantenere. I vecchi comportamenti non si puniscono, si smette di premiarli: la punizione farebbe nascondere gli errori.
\item \textbf{Rinforzo vicario (Bandura):} il messaggio è pubblico, demo day e storia del mese mostrano chi è stato premiato e perché; si impara osservando.
\item \textbf{Attribuzione (Kelley, Weiner):} \guillemotleft{}hai documentato la procedura\guillemotright{} e non \guillemotleft{}sei bravo\guillemotright{}: una causa che la persona controlla si ripete. Se un intero team è silenzioso, la causa è probabilmente il contesto, non le persone.
\item \textbf{Valutazione cognitiva (Deci e Ryan):} premi materiali troppo presto spengono la motivazione intrinseca; per questo prima solo messaggi, punti da M6, premi non promessi in anticipo.
\end{itemize}

\subsection{Limiti}\label{sec:13.5}
La recognition è una leva culturale: i messaggi non rappresentano tutti i comportamenti reali e non sostituiscono log e controlli tecnici. Fiscalità dei premi, privacy e coinvolgimento delle rappresentanze nei 5 paesi sono da verificare con HR e ufficio legale.

\subsection{Compensi per equità e persone}\label{sec:13.6}
Chiediamo a circa 100 persone di fare di più (utenti chiave, Champions, referenti, membri delle strutture temporanee). Per la teoria dell'equità, chi dà di più deve ricevere di più.
\begin{itemize}
\item \textbf{Indennità di ruolo} temporanea (5-10\% del costo annuo) finché dura l'incarico.
\item \textbf{Premio una tantum} (300-500 €) alle persone della linea quando il go-live supera le soglie senza disservizi.
\item \textbf{Progressione stabile} (+5\%) solo per i 20-30 che cambiano davvero ruolo. Nessun aumento generale.
\item \textbf{≈8 assunzioni stabili} dove mancano competenze che servono anche dopo il progetto (3 ERP, 2 dati e AI, 2 sicurezza, 1 HR), tra M3 e M9; \textbf{≈8 sostituti temporanei} per coprire il tempo tolto dal progetto, che finiscono con esso. Prima sempre la mobilità interna.
\end{itemize}

\section{Roadmap dei 24 mesi}\label{sec:14}

{\footnotesize\renewcommand{\arraystretch}{1.25}
\begin{xltabular}{\linewidth}{@{}>{\raggedright\arraybackslash\hsize=0.553\hsize\hspace{0pt}}X>{\raggedright\arraybackslash\hsize=1.000\hsize\hspace{0pt}}X>{\raggedright\arraybackslash\hsize=1.809\hsize\hspace{0pt}}X>{\raggedright\arraybackslash\hsize=0.952\hsize\hspace{0pt}}X>{\raggedright\arraybackslash\hsize=0.838\hsize\hspace{0pt}}X>{\raggedright\arraybackslash\hsize=1.181\hsize\hspace{0pt}}X>{\raggedright\arraybackslash\hsize=0.667\hsize\hspace{0pt}}X@{}}
\toprule
\textbf{Periodo} & \textbf{Ricerca-azione} & \textbf{Programma / ERP} & \textbf{GenAI} & \textbf{Automazioni} & \textbf{Persone e cultura} & \textbf{Misure} \\
\midrule
\endhead
\textbf{M1} & Diagnosi → Analisi & Governance, Change \& Risk Office, referenti; KPMG prepara piano e analisi cyber; soglie di go/no-go fissate; linea pilota scelta (provvisoria) & Contratto firmato & — & Annuncio del CEO (perché); trivia; workshop sui 5 comportamenti & \textbf{T0} \\
\textbf{M2} & Feedback → Azione & Avvio del pilota; utenti chiave scelti; workshop sui processi & Ambiente configurato & Account e formazione avanzata dei 32 Champions & Scheda della persona, plenaria, colloqui su richiesta & pulse mensile \\
\textbf{M3-M5} & Azione & Configurazione a cicli brevi, migrazione dati, profili di accesso & — & M3 classificazione del rischio e registro; M4 prime automazioni e call mensili & Recognition solo messaggi; Cyber essentials; sandbox; community; assunzioni da M3 &  \\
\textbf{M4-M6} & Azione & M6 prova generale → \textbf{go/no-go del go-live} & Diffusione e formazione nelle 8 funzioni & M6 revisione e perimetro & M6 punti e catalogo attivi &  \\
\textbf{M7-M10} & Azione → Valutazione; il ciclo si riapre a M10 & M7 go-live (vecchio sistema per 4 settimane); M8 KPMG riduce; M9 ERP core team, vecchio sistema spento; M10 \textbf{go/no-go dell'estensione}, KPMG esce dal pilota & Consolidamento & M10 le automazioni efficaci entrano nell'ERP & Train-the-trainer; premio al go-live se le soglie tengono & \textbf{T1} (M10, con trivia ripetuto) \\
\textbf{M11-M22} & Azione → Valutazione (M16) & 4 ondate (M11, M13, M16, M18), max 2 linee insieme; KPMG in governance 1 giornata al mese & Monitoraggio di uso, qualità, costi & Integrazione nell'ERP a ogni ondata & Premio a fine ondata se le soglie tengono & \textbf{T2} (M16) \\
\textbf{M23-M24} & Valutazione & Chiusura: valutazione finale, bilancio dei benefici, strutture temporanee sciolte & Gestione ordinaria (IT) & Gestione ordinaria & Fine delle indennità di ruolo & \textbf{T3} \\
\bottomrule
\end{xltabular}}

\textbf{Le due decisioni che contano}
\begin{itemize}
\item \textbf{M6 · go/no-go del go-live:} si accende il nuovo ERP solo se \textbf{test, dati e competenze} degli utenti chiave superano le soglie fissate a M1.
\item \textbf{M10 · go/no-go dell'estensione:} si passa alle altre linee, e si sblocca la \textbf{tranche 2}, solo se il pilota ha tenuto \textbf{continuità, errori e commitment} entro le soglie.
\end{itemize}

\section{Costi, tranche e benefici}\label{sec:15}

\subsection{Le voci di costo (24 mesi, valori centrali, M€)}\label{sec:15.1}

{\small\renewcommand{\arraystretch}{1.25}
\begin{xltabular}{\linewidth}{@{}>{\raggedright\arraybackslash\hsize=0.258\hsize\hspace{0pt}}X>{\raggedright\arraybackslash\hsize=2.453\hsize\hspace{0pt}}X>{\raggedright\arraybackslash\hsize=0.422\hsize\hspace{0pt}}X>{\raggedright\arraybackslash\hsize=1.315\hsize\hspace{0pt}}X>{\raggedright\arraybackslash\hsize=0.552\hsize\hspace{0pt}}X@{}}
\toprule
\textbf{\#} & \textbf{Voce} & \textbf{M€} & \textbf{Area} & \textbf{Tipo} \\
\midrule
\endhead
1 & KPMG sulle 4 ondate & 2,75 & Tecnologia e KPMG & esborso \\
2 & Sicurezza cyber: strumenti e servizi nuovi & 1,50 & Tecnologia e KPMG & esborso \\
3 & Infrastruttura cloud e piattaforma dati/AI & 1,50 & Tecnologia e KPMG & esborso \\
4 & Assistente AI per 850, con consumi extra & 1,40 & Tecnologia e KPMG & esborso \\
5 & Change \& Risk Office e ≈40 referenti & 1,25 & Persone e governance & tempo interno \\
6 & Licenze ERP (da 65 a ≈500 utenti) & 1,10 & Tecnologia e KPMG & esborso \\
7 & Integrazioni e migrazione dei dati & 1,00 & Tecnologia e KPMG & esborso \\
8 & KPMG sul pilota (≈675 giornate × ≈1.300 €) & 0,90 & Tecnologia e KPMG & esborso \\
9 & Compensi per equità & 0,90 & Persone e governance & esborso \\
10 & Assunzioni stabili (≈8 profili) & 0,90 & Persone e governance & esborso \\
11 & Recognition (≈0,35 l'anno da M3) & 0,65 & Persone e governance & esborso \\
12 & Sostituti temporanei (≈8) & 0,60 & Persone e governance & esborso \\
13 & ERP core team (M9-M22) & 0,50 & Tecnologia e KPMG & tempo interno \\
14 & Assistenza dopo ogni go-live (hypercare) & 0,50 & Tecnologia e KPMG & esborso \\
15 & Piattaforma di automazione & 0,45 & Tecnologia e KPMG & esborso \\
16 & Tempo protetto dei 32 Champions & 0,30 & Persone e governance & tempo interno \\
17 & Tempo dei 18 utenti chiave del pilota & 0,25 & Persone e governance & tempo interno \\
18 & Formazione AI (ore interne ed esterna) & 0,22 & Persone e governance & misto \\
19 & KPMG nella governance e controlli sulle automazioni & 0,22 & Persone e governance & esborso \\
20 & Cyber essentials e sandbox ERP & 0,10 & Persone e governance & misto \\
21 & Licenze low-code dei Champions & 0,02 & Persone e governance & esborso \\
 & \textbf{Subtotale} & \textbf{≈17,0} &  &  \\
 & Riserva per imprevisti (20\%) & ≈3,4 &  &  \\
 & \textbf{Totale pianificato} & \textbf{≈20,4} & \textbf{41\% dei 50 M€} &  \\
\bottomrule
\end{xltabular}}

\begin{itemize}
\item \textbf{Tecnologia e KPMG:} ≈11,6 M€ (23\% dei 50).
\item \textbf{Persone e governance del cambiamento:} ≈5,4 M€ (11\%), di cui ≈2,0 M€ di tempo interno, non soldi che escono.
\item \textbf{Nel tempo:} anno 1 ≈6,9 M€, anno 2 ≈10,1 M€ (cumulato ≈17,0 senza riserva).
\end{itemize}

\subsection{Le tranche}\label{sec:15.2}

{\small\renewcommand{\arraystretch}{1.25}
\begin{xltabular}{\linewidth}{@{}>{\raggedright\arraybackslash\hsize=1.375\hsize\hspace{0pt}}X>{\raggedright\arraybackslash\hsize=0.442\hsize\hspace{0pt}}X>{\raggedright\arraybackslash\hsize=0.442\hsize\hspace{0pt}}X>{\raggedright\arraybackslash\hsize=1.742\hsize\hspace{0pt}}X@{}}
\toprule
\textbf{Tranche} & \textbf{Periodo} & \textbf{Importo} & \textbf{Quando si sblocca} \\
\midrule
\endhead
Tranche 1 · pilota e avvio & M1-M10 & ≈7 M€ & all'avvio \\
Tranche 2 · ondate e chiusura & M11-M24 & ≈13 M€ & solo dopo il go/no-go di M10 \\
Non impegnato & — & ≈30 M€ & resta al CdA per nuove esigenze approvate \\
\bottomrule
\end{xltabular}}

Entrambe le tranche includono la loro quota di riserva. Il no-go ferma la spesa: è il modo per spendere solo quando le soglie dicono che conviene.

\subsection{Da dove vengono le stime}\label{sec:15.3}
\begin{itemize}
\item \textbf{ERP:} le fonti indicano 1-5 milioni di dollari per un'implementazione in un'azienda di queste dimensioni e in più paesi; la nostra stima è prudente. Abbiamo aggiunto le voci che di solito si dimenticano: migrazione dei dati (spesso il 25-30\% dell'implementazione, contro il 10\% pianificato), integrazioni, hypercare (3-10\%). Le fonti avvertono che i costi di implementazione si sottostimano spesso del 50-100\%: per questo la riserva è al 20\%.
\item \textbf{KPMG:} le tariffe non sono pubbliche; usiamo ≈1.300 €/giorno come media del team (in Europa, per progetti ERP di fascia media, 1.200-1.800 €/giorno; le Big Four in Germania stanno sopra). Pilota: manager 90 giornate, consulenti 540, cyber 45.
\item \textbf{AI:} scenario base 600 utenti × 25 €/mese, 200 × 50, 50 × 150 (≈0,39 M€/anno) più consumi delle automazioni (≈0,036 M€/anno) e configurazione (25.000 €); consumi extra nello scenario medio +0,25 M€/anno. Uno scenario alto (+0,9 M€/anno) è un rischio coperto dalla riserva.
\item \textbf{Sicurezza e infrastruttura:} ordini di grandezza da benchmark (spesa cyber ≈10-11\% della spesa IT), da affinare con CISO e CIO.
\end{itemize}

\subsection{Benefici attesi (stima prudente)}\label{sec:15.4}

{\small\renewcommand{\arraystretch}{1.25}
\begin{xltabular}{\linewidth}{@{}>{\raggedright\arraybackslash\hsize=0.612\hsize\hspace{0pt}}X>{\raggedright\arraybackslash\hsize=2.545\hsize\hspace{0pt}}X>{\raggedright\arraybackslash\hsize=0.454\hsize\hspace{0pt}}X>{\raggedright\arraybackslash\hsize=0.389\hsize\hspace{0pt}}X@{}}
\toprule
\textbf{Fonte} & \textbf{Calcolo} & \textbf{Ore l'anno} & \textbf{M€ l'anno} \\
\midrule
\endhead
ERP & ≈500 utenti × 5\% del tempo, dopo la stabilizzazione & ≈40.000 & ≈1,2 \\
Assistente AI & ≈500 utenti d'ufficio × 2\% del tempo (le indagini indicano 1-2\% sulla media, 3-5\% per chi lo usa) & ≈16.000 & ≈0,5 \\
Automazioni & 32 Champions × ≈5 automazioni × ≈8 ore risparmiate al mese & ≈15.000 & ≈0,45 \\
\textbf{Totale a regime} &  & \textbf{≈70.000} & \textbf{≈2,1} \\
\bottomrule
\end{xltabular}}

\begin{itemize}
\item 70.000 ore ≈ 45 persone a tempo pieno, ≈5\% della forza lavoro; pieno dal primo anno dopo la chiusura.
\item \textbf{Calo iniziale dichiarato:} −15\% di produttività per 2 mesi sulle ≈60 persone del pilota (≈0,07 M€).
\item \textbf{Le ore non diventano tagli:} si reinvestono e alimentano la mobilità interna verso i ruoli nuovi.
\item \textbf{Che cosa non ripagano:} le ore liberate coprono i costi che restano dopo il progetto (licenze, profili assunti, progressioni), ma non ripagano da sole l'investimento. Il resto del valore sta in ciò che non quantifichiamo: dati affidabili, meno errori, rischio cyber sotto controllo, possibilità di crescere senza assumere in proporzione. Lo dichiariamo, perché un piano credibile dichiara i suoi limiti.
\end{itemize}

\section{Valutazione: KPI, commitment e benessere}\label{sec:16}

\subsection{Principi}\label{sec:16.1}
\begin{itemize}
\item \textbf{Due domande:} il cambiamento funziona (risultati)? Come lo vivono le persone (commitment, benessere, accoglienza della recognition)?
\item \textbf{Chi vede che cosa:} la governance vede solo \textbf{dati aggregati} (almeno 10 persone per gruppo). I KPI individuali li vedono solo la persona e HR, come guida per il supporto. \textbf{Nessun KPI valuta la prestazione delle persone.}
\item \textbf{Soglie:} quelle indicate sono \textbf{esempi, da tarare su T0 con la governance a M1}, sulla situazione di partenza. Il pilota è il benchmark ottimistico; le ondate hanno soglie più prudenti.
\item \textbf{Calendario:} \textbf{T0} (M1), \textbf{T1} (M10, fine pilota, con trivia ripetuto), \textbf{T2} (M16, metà ondate), \textbf{T3} (M23-M24, chiusura), più un \textbf{pulse mensile} breve.
\item \textbf{Fonti dei dati:} registri dell'ERP, console aggregata dell'AI, registro delle automazioni, piattaforma di recognition, registro della formazione, trivia, questionari anonimi.
\item Se un indicatore esce dalle soglie, la governance \textbf{corregge prima di andare avanti}: è la ricerca-azione applicata alla gestione del progetto. Se un gruppo resta sotto soglia, la causa si cerca prima nel contesto (formazione, carico, strumenti).
\end{itemize}

\subsection{Commitment al cambiamento (Herscovitch e Meyer)}\label{sec:16.2}
Il questionario \emph{Commitment to Change} (18 item, 6 per componente) distingue tre forme di adesione:
\begin{itemize}
\item \textbf{\guillemotleft{}Voglio\guillemotright{}} (affettivo): la persona ci crede e sostiene il cambiamento.
\item \textbf{\guillemotleft{}Devo\guillemotright{}} (di continuità): opporsi o andarsene costerebbe troppo.
\item \textbf{\guillemotleft{}È giusto\guillemotright{}} (normativo): si sente in dovere di aderire.
\end{itemize}

Gli autori collegano queste forme a un continuum di comportamenti: resistenza attiva → resistenza passiva → adeguamento → cooperazione → \emph{championing}. Solo \guillemotleft{}voglio\guillemotright{} (e in parte \guillemotleft{}è giusto\guillemotright{}) porta a cooperazione e championing; \guillemotleft{}devo\guillemotright{} porta all'adeguamento. \textbf{L'andamento del commitment dice quindi se le persone sostengono il cambiamento o si stanno solo adeguando.} È anche un controllo della recognition: se funziona deve crescere \guillemotleft{}voglio\guillemotright{}; se cresce solo \guillemotleft{}devo\guillemotright{}, il rinforzo produce adeguamento, non adesione.

Il questionario non sostituisce il trivia (che assegna i ruoli): è una misura indipendente, anonima e aggregata, per verificarne la tenuta nel tempo.

\subsection{Benessere e clima}\label{sec:16.3}

{\small\renewcommand{\arraystretch}{1.25}
\begin{xltabular}{\linewidth}{@{}>{\raggedright\arraybackslash\hsize=0.606\hsize\hspace{0pt}}X>{\raggedright\arraybackslash\hsize=1.153\hsize\hspace{0pt}}X>{\raggedright\arraybackslash\hsize=0.290\hsize\hspace{0pt}}X>{\raggedright\arraybackslash\hsize=1.950\hsize\hspace{0pt}}X@{}}
\toprule
\textbf{Strumento} & \textbf{Che cosa misura} & \textbf{Item} & \textbf{Funzione} \\
\midrule
\endhead
\textbf{WHO-5 Well-Being Index} & Benessere psicologico generale & 5 & Pulse mensile: segnala presto un peggioramento \\
\textbf{HSE Management Standards Indicator Tool} & Stress lavoro-correlato e sue cause organizzative: richieste (carico), controllo, supporto di capi e colleghi, relazioni, ruolo, cambiamento & 35 & Base della metodologia INAIL per lo stress lavoro-correlato (obbligatoria in Italia, D.Lgs. 81/2008); misura le cause per reparto, non la persona. Supporto e relazioni ci dicono anche se le persone si sentono sicure nel segnalare errori. \\
\textbf{Copenhagen Burnout Inventory}, scala lavoro & Esaurimento legato al lavoro & 7 & Coglie il sovraccarico della transizione, soprattutto nei ruoli extra \\
\bottomrule
\end{xltabular}}

\begin{itemize}
\item \textbf{Cadenza:} batteria completa a T0-T3 (≈65 item, circa 12-15 minuti con il Commitment to Change); ogni mese WHO-5 + tre domande rapide (carico, supporto, recognition) + un campo libero.
\item \textbf{Feedback anonimo libero} sempre disponibile.
\item \textbf{Garanzie} (dati sulla salute = categoria particolare, art. 9 GDPR): partecipazione volontaria, anonima, solo aggregati con almeno 10 risposte per incrocio; gestione di HR con il medico competente o un ente terzo, mai dei responsabili diretti.
\item \textbf{Soglia:} nessun peggioramento rispetto a T0; la soglia di rischio che fa scattare un intervento è \textbf{da concordare} con il medico competente sui valori di riferimento degli strumenti. Superata la soglia, scattano \textbf{interventi organizzativi} (carichi, supporto, tempi), mai sul singolo; sportello di ascolto volontario.
\item Indicatori ``sentinella'' a costo zero: assenze, turnover, infortuni, segnalazioni del medico competente.
\end{itemize}

\subsection{I KPI del team di governance (10 per fase)}\label{sec:16.4}
In ogni fase: ≈3 di continuità operativa, 2-3 di rischio tecnologico, ≈3 su persone e conoscenze, 1-2 di budget. Formato: \emph{KPI — metrica · soglia · frequenza · fonte}.

\textbf{Fase 1 · Preparazione (M1-M6): arrivare pronti al go-live}
\begin{enumerate}
\item Configurazione nei tempi — \% cicli chiusi come da piano · ≥90\% · mensile · ERP core team, KPMG
\item Dati pronti alla migrazione — \% dati puliti e validati · ≥98\% entro M6 · mensile · ERP core team
\item Champions delle automazioni attivi — formati e con account sui ≈32 · 100\% entro M4 · mensile · registro formazione
\item Incidenti cyber gravi · 0 · mensile · CISO
\item Cyber essentials — \% dipendenti formati · 100\% entro M5 · mensile · registro formazione
\item AI dentro i profili — \% utenti attivi sull'assistente aziendale · in crescita ogni mese · console AI
\item Trivia e restituzione — \% con trivia completato e scheda aperta · ≥95\% entro M2 · sito, HR
\item Benessere e stress — pulse WHO-5 e HSE rispetto a T0 · nessun peggioramento · mensile
\item Formazione nel tempo protetto — \% ore svolte nelle ore dedicate · ≥90\% · mensile
\item Spesa rispetto alla tranche 1 · ≤100\% · mensile · CFO
\end{enumerate}

\textbf{Fase 2 · Go-live e stabilizzazione (M7-M10): il pilota regge?}
\begin{enumerate}
\item Consegne puntuali della linea — \% ordini in tempo rispetto a T0 · ≥95\% di T0 · mensile · ERP
\item Errori nelle transazioni — errori ogni 1.000 transazioni · di nuovo a T0 entro 8 settimane dal go-live · settimanale · ERP
\item Transazioni nel nuovo ERP — \% a sistema (strumenti paralleli in calo) · ≥90\% entro M9 · mensile · ERP
\item Incidenti cyber gravi · 0
\item Profili di accesso rivisti prima del go-live · 100\% · CISO, KPMG
\item Automazioni con istruzioni e sostituto · 100\% · registro
\item Benessere e stress · nessun peggioramento da T0
\item Copertura della recognition — \% che danno e ricevono a trimestre · ≥70\%, divario presenza/remoto <10 punti · piattaforma
\item Dipendenza da KPMG — ore KPMG nel team ERP · in calo da M8 · Change \& Risk Office
\item Spesa rispetto alla tranche 1 · ≤100\%
\end{enumerate}

\textbf{Fase 3 · Estensione (M11-M22): il modello si ripete?}
\begin{enumerate}
\item Consegne puntuali delle linee in ondata · ≥92\% di T0 · mensile
\item Transazioni nel nuovo ERP delle linee in ondata · ≥90\% entro 3 mesi dal go-live
\item Ondate formate internamente — \% formazione fatta da personale NexaCore · ≥80\% · a ogni ondata
\item Incidenti cyber gravi · 0
\item AI fuori dagli strumenti aziendali — quota che dichiara di usare strumenti personali · in calo ogni trimestre · pulse anonimo
\item Automazioni integrate nell'ERP · ≥1 per linea · a ogni ondata · registro
\item Commitment al cambiamento — quota \guillemotleft{}voglio\guillemotright{} rispetto a \guillemotleft{}devo\guillemotright{} · in crescita da T0 · T1-T2
\item Movimento tra categorie — \% Fence sitters diventati Champions · in crescita da T0 · trivia ripetuto
\item Turnover volontario (totale e ruoli extra) · ≤ T0 · trimestrale · HR
\item Spesa rispetto alla tranche 2 · ≤100\%
\end{enumerate}

\textbf{Fase 4 · Chiusura (M23-M24): il cambiamento è diventato ordinario?}
\begin{enumerate}
\item ERP in tutta l'azienda — \% transazioni di tutte le linee a sistema · ≥95\%
\item Tempo del processo ordine-incasso rispetto a T0 · ≤ T0
\item Ore liberate con reinvestimento dichiarato · 100\%
\item Incidenti cyber gravi · 0
\item Clic nelle simulazioni di phishing · ≤5\% · trimestrale · CISO
\item Valore delle ore risparmiate / costo dei token · > 1 per ogni automazione · registro
\item Commitment: quota \guillemotleft{}voglio\guillemotright{} a T3 rispetto a T0 · in crescita
\item Benessere e burnout dei ruoli extra (WHO-5, HSE, CBI) · nessun peggioramento da T0
\item Presidio passato alla gestione ordinaria — ruoli, procedure e sostituti formalizzati · 100\%
\item Spesa totale rispetto al budget · ≤100\%
\end{enumerate}

I KPI di continuità e di rischio alimentano i go/no-go di M6, M10 e di ogni ondata.

\subsection{KPI per ruolo nel cambiamento}\label{sec:16.5}

{\small\renewcommand{\arraystretch}{1.25}
\begin{xltabular}{\linewidth}{@{}>{\raggedright\arraybackslash\hsize=0.329\hsize\hspace{0pt}}X>{\raggedright\arraybackslash\hsize=1.671\hsize\hspace{0pt}}X@{}}
\toprule
\textbf{Ruolo} & \textbf{KPI principali} \\
\midrule
\endhead
Utente chiave ERP & Colleghi formati e autonomi (tutta la funzione entro il go-live); richieste risolte senza KPMG (≥50\% a M8, ≥80\% a M10); manuale della linea completo entro M10 \\
Champion delle automazioni & Automazioni attive (≈5 a fine anno 1); 100\% con istruzioni, sostituto e costo stimato; ≈8 ore risparmiate al mese per automazione \\
Referente di funzione e sede & Copertura recognition della funzione ≥70\% a trimestre; almeno una segnalazione al mese al Change \& Risk Office \\
Change \& Risk Office & Resoconto ``ci avete detto / abbiamo cambiato'' ogni 2 mesi; report mensile alla governance (12 su 12) \\
ERP core team & Rilasci senza blocchi ≥95\%; ondate formate internamente ≥80\% \\
Process owner & Tempo del processo ≤ T0 entro 3 mesi dal go-live \\
\bottomrule
\end{xltabular}}

\subsection{KPI per persona (sempre 3)}\label{sec:16.6}
Ogni persona ha \textbf{3 KPI}, che misurano \textbf{solo come attraversa il cambiamento} (adozione, apprendimento, partecipazione, condivisione), non il rendimento nel lavoro ordinario. Ogni KPI ha una variante \emph{prima} e \emph{dopo} il momento in cui cambia il lavoro (go-live della propria linea; per chi non è nell'ERP, M6 per l'AI e M10 per il trivia ripetuto). Alcuni KPI misurano l'azienda, non la persona: se gli accordi o la formazione promessa mancano, il KPI è dell'azienda.

\textbf{KPI 1 · Adozione (per tutti)}
\begin{itemize}
\item Team ERP: preparazione (100\% di formazione e prove in sandbox entro il go-live) → adozione (≥90\% delle proprie attività registrate nell'ERP entro 3 mesi).
\item Non ERP: formazione di base sull'AI (4 ore entro M5) → uso dell'assistente almeno 3 settimane su 4 da M6.
\end{itemize}

\textbf{KPI 2 e 3 (per persona)}

{\small\renewcommand{\arraystretch}{1.25}
\begin{xltabular}{\linewidth}{@{}>{\raggedright\arraybackslash\hsize=0.566\hsize\hspace{0pt}}X>{\raggedright\arraybackslash\hsize=1.141\hsize\hspace{0pt}}X>{\raggedright\arraybackslash\hsize=1.293\hsize\hspace{0pt}}X@{}}
\toprule
\textbf{Persona} & \textbf{KPI 2 · la persona} & \textbf{KPI 3 · il ruolo o la persona nel ruolo} \\
\midrule
\endhead
111 Pioniere & AI entro le regole (0 usi fuori profilo) & ERP: colleghi formati → richieste risolte senza KPMG. Non ERP: formazione avanzata → automazioni attive \\
110 Talento in bilico & Accordi su carico e riconoscimento rispettati (100\%) → equità percepita in crescita da T0 & come il Pioniere \\
011 Volenteroso in crescita & ERP: ≥10 esercizi in sandbox → compiti in autonomia. Non ERP: formazione AI → autoefficacia sopra 50 a T1 & come il Pioniere \\
101 Esecutore in cerca di senso & Workshop di aggiornamento (≥80\%) & Una prova scelta con il responsabile / una prima iniziativa → motivazione in crescita da T0 \\
010 Contributore da sostenere & Formazione nel tempo protetto (100\% delle ore) & Prove in coppia (ERP ≥3 prima del go-live; AI ≥2 entro M8) → compiti in autonomia \\
001 Esploratore prudente & Prove guidate / primo compito guidato → compiti con verifica dell'output & Affiancamento svolto → autoefficacia (o fiducia nel controllare l'AI) in crescita \\
100 Esperto in cerca di garanzie & Critiche che cambiano il processo (≥1 a trimestre) & Conoscenza trasferita (≥2 procedure o casi limite) → strumenti paralleli abbandonati \\
000 Professionista in cerca di stabilità & Prova circoscritta → compiti in autonomia (o con aiuto) & Impegni dell'azienda rispettati (accordi con HR mantenuti, 100\%) \\
\bottomrule
\end{xltabular}}

\section{Matrice dei timori del Board}\label{sec:17}

{\small\renewcommand{\arraystretch}{1.25}
\begin{xltabular}{\linewidth}{@{}>{\raggedright\arraybackslash\hsize=0.667\hsize\hspace{0pt}}X>{\raggedright\arraybackslash\hsize=1.736\hsize\hspace{0pt}}X>{\raggedright\arraybackslash\hsize=0.716\hsize\hspace{0pt}}X>{\raggedright\arraybackslash\hsize=1.263\hsize\hspace{0pt}}X>{\raggedright\arraybackslash\hsize=0.618\hsize\hspace{0pt}}X@{}}
\toprule
\textbf{Timore} & \textbf{Leve principali} & \textbf{Chi presidia} & \textbf{KPI} & \textbf{Soglia (da tarare su T0)} \\
\midrule
\endhead
\textbf{Resistenza al cambiamento} & Trivia e ruoli; Champions formatori; Cinici consultati; cerchio d'oro; recognition; nessun taglio legato alle automazioni & Governance, HR, referenti & Commitment \guillemotleft{}voglio\guillemotright{} rispetto a \guillemotleft{}devo\guillemotright{}; Fence sitters diventati Champions; turnover volontario & in crescita da T0; in crescita da T0; ≤ T0 \\
\textbf{Perdita di conoscenze chiave} & ERP core team; utenti chiave formatori; manuale della linea; sostituto per ogni automazione; mentoring senior-junior; knowledge base & CIO, ERP core team, Change \& Risk Office & Formazione delle ondate fatta da personale interno; automazioni con istruzioni e sostituto; richieste risolte senza KPMG & ≥80\%; 100\%; ≥80\% a M10 \\
\textbf{Errori nella transizione} & Una sola linea; vecchio sistema per 4 settimane; sandbox; competenza come soglia di go-live; revisione degli errori senza colpe; max 2 linee in ondata & Responsabile di linea, process owner, governance & Consegne puntuali; errori ogni 1.000 transazioni; transazioni nel nuovo ERP & ≥95\% di T0; a T0 entro 8 settimane; ≥90\% entro M9 \\
\textbf{Nuovi rischi cyber} & CISO in governance; esperto cyber KPMG; accessi rivisti prima di ogni go-live; Cyber essentials per 850; investimento in sicurezza (≈1,5 M€) & CISO & Incidenti gravi; profili di accesso controllati; clic nelle simulazioni di phishing & 0; 100\%; ≤5\% \\
\textbf{AI non governata} & Contratto enterprise; tre profili; classificazione del rischio; registro; recognition per chi registra; tetti di spesa & Governance, Change \& Risk Office, IT & Uso dichiarato di strumenti personali; automazioni nel registro con sostituto; valore delle ore rispetto al costo & in calo ogni trimestre; 100\%; > costo \\
\bottomrule
\end{xltabular}}

\textbf{Benefici e rischi residui per tecnologia}

{\small\renewcommand{\arraystretch}{1.25}
\begin{xltabular}{\linewidth}{@{}>{\raggedright\arraybackslash\hsize=0.401\hsize\hspace{0pt}}X>{\raggedright\arraybackslash\hsize=1.253\hsize\hspace{0pt}}X>{\raggedright\arraybackslash\hsize=1.346\hsize\hspace{0pt}}X@{}}
\toprule
\textbf{Timore} & \textbf{Pilota ERP} & \textbf{Automazioni} \\
\midrule
\endhead
Resistenza & \textbf{+} parte da chi è favorevole. \textbf{−} risultati migliori del reale. \emph{Mitigazione:} benchmark ottimistico, Cinici coinvolti & \textbf{+} il cambiamento parte dai colleghi. \textbf{−} paura per il posto. \emph{Mitigazione:} niente tagli, ore reinvestite \\
Conoscenze & \textbf{+} competenze trasferite, manuale. \textbf{−} dipendenza da KPMG. \emph{Mitigazione:} ritiro graduale legato all'autonomia & \textbf{+} il processo viene scritto. \textbf{−} lo conosce solo l'autore. \emph{Mitigazione:} sostituto obbligatorio \\
Errori & \textbf{+} confinati a una linea. \textbf{−} doppio lavoro, disservizio. \emph{Mitigazione:} sistema in parallelo, piano di ritorno & \textbf{+} meno errori manuali. \textbf{−} errori replicati in serie. \emph{Mitigazione:} controllo umano, classificazione del rischio \\
Cyber & \textbf{+} accessi progettati con KPMG. \textbf{−} esposizione del cloud. \emph{Mitigazione:} revisione degli accessi prima del go-live & \textbf{+} l'AI nascosta viene alla luce. \textbf{−} più dati nell'AI. \emph{Mitigazione:} solo strumento enterprise, dati personali e pagamenti esclusi \\
AI non governata & Ospita le prime automazioni sotto controllo & \textbf{+} registro e strumento unico. \textbf{−} automazioni fuori registro. \emph{Mitigazione:} il registro dà riconoscimento \\
\bottomrule
\end{xltabular}}

\textbf{Rispetto ai tre obiettivi del caso}
\begin{itemize}
\item \textbf{Continuità operativa:} una linea alla volta, il team più pronto, rete di sicurezza, soglie di go/no-go, max due linee in transizione.
\item \textbf{Controllo dei rischi tecnologici:} CISO e KPMG in governance, classificazione del rischio e registro, accessi rivisti, contratto AI enterprise, tranche legate alle soglie.
\item \textbf{Valorizzazione delle persone:} ruoli nel cambiamento per i Champions, formazione in orario, compensi per equità, recognition tra pari, over 50 valorizzati nel mentoring, ascolto dei Cinici, benessere monitorato.
\end{itemize}

\section{Rischi propri del piano}\label{sec:18}

{\small\renewcommand{\arraystretch}{1.25}
\begin{xltabular}{\linewidth}{@{}>{\raggedright\arraybackslash\hsize=0.648\hsize\hspace{0pt}}X>{\raggedright\arraybackslash\hsize=1.352\hsize\hspace{0pt}}X@{}}
\toprule
\textbf{Rischio} & \textbf{Mitigazione} \\
\midrule
\endhead
Champions sovraccarichi & Tempo protetto, sostituti temporanei, pulse sul carico, indennità di ruolo \\
Sforamenti di tempi e budget, frequenti nei progetti ERP & Go/no-go su soglie fissate a M1, tranche, riserva del 20\%, controllo mensile della spesa \\
Più Champions sulla carta che nei fatti (trivia nominativo) & Restituzione, osservazione dei comportamenti, trivia ripetuto, incarichi volontari \\
Dipendenza da KPMG & Uscita programmata dal pilota, ERP core team, utenti chiave formatori \\
Struttura soft che resta sulla carta & Tempo formalizzato e obiettivi nella valutazione dei referenti \\
Premi che spiazzano la motivazione intrinseca & Avvio senza punti, premi non promessi in anticipo, nessun punto per il volume d'uso \\
Abuso dei punti & Tetti, divieto di auto-assegnazione, avvisi sugli scambi reciproci, distribuzione pubblicata \\
Consumi dell'AI oltre il previsto & Tetti per profilo, costo stimato nel registro, scenario alto coperto dalla riserva \\
Uso improprio dei dati del trivia & Accesso limitato a persona, HR e responsabile del pilota; mai per valutare; aggregati con almeno 10 persone; informativa e verifica legale \\
Vincoli legali nei 5 paesi (privacy, fiscalità dei premi, rappresentanze) & Verifica con HR e ufficio legale prima dell'avvio \\
\bottomrule
\end{xltabular}}

\section{Punti aperti}\label{sec:19}

Li dichiariamo perché sono decisioni che, in un progetto reale, spetterebbero alla governance di NexaCore o richiederebbero dati che non abbiamo.
\begin{enumerate}
\item \textbf{Numero di linee di prodotto.} Il piano assume ≈5 linee (pilota + 4 ondate) su 4 settori di clienti. Se le linee fossero 4, le ondate sarebbero 3: o il programma si accorcia (e si ricalcolano i costi KPMG), o le ondate si distanziano nei 24 mesi.
\item \textbf{Soglie:} tutte da tarare su T0 con la governance a M1.
\item \textbf{Costi di sicurezza e infrastruttura:} ordini di grandezza da affinare con CISO e CIO.
\item \textbf{Costo orario:} usiamo il dato italiano; quello pesato sui 5 paesi è probabilmente più alto.
\item \textbf{KPMG sulle ondate:} si potrebbe ridurre grazie all'ERP core team.
\item \textbf{Conformità del trivia nominativo:} verifica legale nei 5 paesi (\hyperref[sec:5.6]{\S5.6}).
\item \textbf{Uno strumento specifico per la sicurezza psicologica} (propensione a segnalare errori): oggi la leggiamo attraverso le scale di supporto e relazioni dell'HSE e il numero di segnalazioni; uno strumento dedicato è da valutare.
\item \textbf{Proposta non decisa:} un piano di azionariato aperto a tutti i dipendenti, alimentato dai punti della recognition.
\end{enumerate}

\section{Domande frequenti}\label{sec:20}

\textbf{Perché un trivia obbligatorio e nominativo? Non rischia di spaventare le persone?}
Serve ad assegnare ruoli concreti (chi diventa utente chiave, chi automatizza, quale linea fa da pilota): con dati anonimi non si potrebbe. Le garanzie sono strette: il profilo lo vedono solo la persona, HR e il responsabile del pilota; non serve mai a valutare; ognuno lo riceve e può contestarlo; gli incarichi restano volontari (\hyperref[sec:5.6]{\S5.6}, \hyperref[sec:7]{\S7}).

\textbf{Perché non riorganizzare l'azienda per processi?}
Cambiare la gerarchia durante tre transizioni tecnologiche aggiungerebbe incertezza e rischio operativo. Usiamo i team di prodotto già esistenti e strutture temporanee che si sciolgono a fine programma (\hyperref[sec:8]{\S8}).

\textbf{Perché il pilota su una linea di prodotto e non su un reparto o una sede?}
Una linea attraversa tutti i reparti, quindi collauda i passaggi tra funzioni dove nascono gli errori. Un reparto sarebbe poco rappresentativo, una sede troppo rischiosa (\hyperref[sec:9.1]{\S9.1}).

\textbf{Scegliere il team con più Champions non falsa i risultati?}
Sì, e lo dichiariamo: i risultati del pilota sono un benchmark ottimistico. In cambio proteggiamo i clienti nella fase più rischiosa. Le ondate hanno soglie più prudenti (\hyperref[sec:9.2]{\S9.2}).

\textbf{Che cosa succede se un go/no-go dà esito negativo?}
Non si procede: a M6 il go-live si rinvia finché test, dati e competenze non sono pronti; a M10 non si estende e la tranche 2 non si sblocca. La governance decide le correzioni e ripete la verifica (\hyperref[sec:14]{\S14}, \hyperref[sec:15.2]{\S15.2}).

\textbf{KPMG non è in conflitto di interessi?}
Il rischio c'è (configura e consiglia sul go/no-go). Lo mitighiamo con persone diverse, soglie scritte prima, verbali che separano parere e decisione, e ruolo senza voto (\hyperref[sec:8.3]{\S8.3}).

\textbf{Che cosa resta dopo l'uscita di KPMG?}
L'ERP core team, gli utenti chiave formati come formatori, il manuale della linea, la knowledge base; a fine programma le attività passano a IT, alle linee e ai profili assunti (\hyperref[sec:8.2]{\S8.2}, \hyperref[sec:9.5]{\S9.5}).

\textbf{Perché non vietare semplicemente l'AI non autorizzata?}
Perché il divieto senza alternativa spinge l'uso nell'ombra. Diamo a tutti uno strumento aziendale sicuro, regole chiare e un vantaggio a registrare le automazioni (\hyperref[sec:10]{\S10}, \hyperref[sec:11]{\S11}, \hyperref[sec:13]{\S13}).

\textbf{Chi decide che cosa si può automatizzare?}
Il team di governance, con la classificazione del rischio, su parere consultivo di KPMG. Dati personali e pagamenti sono esclusi all'inizio (\hyperref[sec:11.2]{\S11.2}).

\textbf{Le automazioni porteranno tagli di personale?}
No: le ore liberate si reinvestono in attività di valore e nella mobilità interna. È una proposta da formalizzare come impegno aziendale (A11, \hyperref[sec:11.2]{\S11.2}).

\textbf{Perché spendere solo 20 dei 50 milioni?}
Perché il piano dimensiona ciò che serve davvero, con una riserva del 20\%. Il resto (≈30 M€) non viene impegnato: resta al CdA per esigenze approvate in seguito. Le tranche fanno spendere solo quando le soglie lo giustificano (\hyperref[sec:15]{\S15}).

\textbf{L'investimento si ripaga?}
Non con le sole ore liberate (≈2,1 M€ l'anno). Il valore principale sta in dati affidabili, meno errori, rischio cyber sotto controllo e capacità di crescere senza assumere in proporzione. Lo diciamo apertamente (\hyperref[sec:15.4]{\S15.4}).

\textbf{Come si capisce se le persone aderiscono davvero o si adeguano?}
Con il commitment al cambiamento di Herscovitch e Meyer: se cresce \guillemotleft{}voglio\guillemotright{} le persone sostengono il cambiamento; se cresce solo \guillemotleft{}devo\guillemotright{} si stanno adeguando (\hyperref[sec:16.2]{\S16.2}).

\textbf{Come si evita che il benessere peggiori durante la transizione?}
Pulse mensile (WHO-5 + carico, supporto, recognition), questionari HSE e burnout a T0-T3, soglia ``nessun peggioramento da T0''; se peggiora si interviene su carichi e supporto, mai sul singolo (\hyperref[sec:16.3]{\S16.3}).

\textbf{Che cosa si sa già di NexaCore e che cosa è stimato?}
Dal caso: dimensioni, sedi, reparti, età, lavoro remoto, 50 M€, timori del Board. Tutto il resto (linee, persone per ruolo, costi, soglie) è stima o assunzione del gruppo (\hyperref[sec:3]{\S3}).

\section{Glossario}\label{sec:21}

\begin{itemize}
\item \textbf{Champions / Fence sitters / Cinici:} le tre categorie di atteggiamento verso il cambiamento in cui il trivia colloca le persone: chi lo guida, chi aspetta prima di esporsi, chi rema contro ma segnala problemi reali (\hyperref[sec:5.3]{\S5.3}).
\item \textbf{Persona:} uno degli 8 profili del trivia, identificato da un codice di tre cifre (autoefficacia, autodeterminazione, equità; 1 alta, 0 bassa).
\item \textbf{Trivia:} questionario a scenari di 12 domande che profila le persone (\hyperref[sec:5]{\S5}).
\item \textbf{Ricerca-azione:} approccio al cambiamento che alterna raccolta di dati, restituzione alle persone, azione e nuova misura (\hyperref[sec:2.2]{\S2.2}).
\item \textbf{Struttura hard / soft:} l'organigramma formale / le relazioni, i flussi e i team che attraversano l'organigramma.
\item \textbf{Team di prodotto:} team interfunzionale già esistente che gestisce una linea di prodotto; unità del pilota e delle ondate.
\item \textbf{Team di governance:} organo temporaneo che decide soglie, go/no-go, classificazione del rischio e budget (\hyperref[sec:8.2]{\S8.2}).
\item \textbf{Change \& Risk Office:} unità temporanea in staff al CEO che coordina il programma (\hyperref[sec:8.2]{\S8.2}).
\item \textbf{Utente chiave:} Champion del team pilota che prova i processi, forma i colleghi e poi le linee delle ondate.
\item \textbf{ERP core team:} squadra interna temporanea (M9-M22) che prende il posto di KPMG nella gestione dell'ERP.
\item \textbf{Process owner:} responsabile del risultato di un processo end-to-end (es. ordine-incasso).
\item \textbf{Go/no-go:} decisione formale, su soglie fissate in anticipo, se procedere a una fase successiva.
\item \textbf{Go-live:} accensione del nuovo ERP per una linea.
\item \textbf{Ondata:} passaggio all'ERP di una linea dopo il pilota.
\item \textbf{Hypercare:} assistenza rafforzata nelle settimane successive a un go-live.
\item \textbf{Sandbox:} ambiente di prova dell'ERP con dati fittizi.
\item \textbf{Train-the-trainer:} formazione che prepara una persona a formare a sua volta i colleghi.
\item \textbf{Classificazione del rischio:} regola della governance che dice che cosa si può automatizzare e con quali approvazioni.
\item \textbf{Registro delle automazioni:} elenco di ogni automazione con processo, dati, autore, sostituto, ore risparmiate e costo.
\item \textbf{Shadow AI:} uso di strumenti di AI non autorizzati, con dati aziendali fuori controllo.
\item \textbf{Token:} unità di testo con cui si misura e si paga l'uso dei modelli linguistici.
\item \textbf{Recognition:} riconoscimento pubblico tra colleghi dei comportamenti utili, con tag e, da M6, punti.
\item \textbf{Tranche:} quota del budget che si sblocca solo dopo un go/no-go positivo.
\item \textbf{T0, T1, T2, T3:} momenti di misura completa (M1, M10, M16, M23-M24).
\item \textbf{Pulse:} breve questionario mensile su benessere, carico, supporto e recognition.
\item \textbf{Commitment al cambiamento:} adesione al cambiamento nelle tre forme \guillemotleft{}voglio\guillemotright{}, \guillemotleft{}devo\guillemotright{}, \guillemotleft{}è giusto\guillemotright{} (\hyperref[sec:16.2]{\S16.2}).
\end{itemize}

\section{Fondamenti teorici: dove li usiamo}\label{sec:22}

{\small\renewcommand{\arraystretch}{1.25}
\begin{xltabular}{\linewidth}{@{}>{\raggedright\arraybackslash\hsize=1.193\hsize\hspace{0pt}}X>{\raggedright\arraybackslash\hsize=0.591\hsize\hspace{0pt}}X>{\raggedright\arraybackslash\hsize=1.841\hsize\hspace{0pt}}X>{\raggedright\arraybackslash\hsize=0.375\hsize\hspace{0pt}}X@{}}
\toprule
\textbf{Teoria} & \textbf{Autori} & \textbf{Scelta che motiva} & \textbf{Dove} \\
\midrule
\endhead
Ricerca-azione: diagnosi, analisi, feedback, azione, valutazione & Lewin & Metodo del piano; fasi della roadmap & \hyperref[sec:2]{\S2}, \hyperref[sec:14]{\S14} \\
Otto passi del cambiamento, vittorie rapide & Kotter & Pilota con il team più pronto & \hyperref[sec:9]{\S9}, \hyperref[sec:12.3]{\S12.3} \\
Campioni, seduti sulla staccionata, cinici & Noble (richiamato nel corso) & Le tre categorie & \hyperref[sec:5.3]{\S5.3} \\
Autoefficacia; rinforzo vicario & Bandura & Asse del trivia; formazione per padronanza; recognition pubblica & \hyperref[sec:5]{\S5}, \hyperref[sec:12]{\S12}, \hyperref[sec:13]{\S13} \\
Autodeterminazione; valutazione cognitiva & Deci e Ryan & Asse del trivia; avvio della recognition senza punti & \hyperref[sec:5]{\S5}, \hyperref[sec:13]{\S13} \\
Equità & Adams & Asse del trivia; formazione in orario, indennità, costi delle persone & \hyperref[sec:5]{\S5}, \hyperref[sec:12]{\S12}, \hyperref[sec:13.6]{\S13.6}, \hyperref[sec:15]{\S15} \\
Bisogni di potere, affiliazione, realizzazione & McClelland & Affinità tra Champions e ruoli & \hyperref[sec:8.4]{\S8.4} \\
Definizione degli obiettivi & Locke & Soglie di go-live specifiche e note in anticipo & \hyperref[sec:12.3]{\S12.3} \\
Aspettativa & Vroom & Formazione legata a risultato e ricompensa & \hyperref[sec:12.3]{\S12.3} \\
Rinforzo & Skinner & Recognition legata al comportamento e immediata & \hyperref[sec:13.4]{\S13.4} \\
Attribuzione & Kelley, Weiner & Messaggi che nominano il comportamento; lettura dei silenzi come problema di contesto & \hyperref[sec:13.4]{\S13.4} \\
Cerchio d'oro & Sinek & Comunicazione del CEO dal perché & \hyperref[sec:13.1]{\S13.1} \\
Commitment al cambiamento & Herscovitch e Meyer & Misura di adesione o adeguamento & \hyperref[sec:16.2]{\S16.2} \\
Bias percettivi, dissonanza cognitiva, ambivalenza & (corso) & Lettura delle reazioni alla restituzione & \hyperref[sec:5.4]{\S5.4}, \hyperref[sec:7.2]{\S7.2} \\
Strutture a matrice e unità di staff & (corso) & Regola anti-ambiguità; Change \& Risk Office & \hyperref[sec:8.2]{\S8.2} \\
Barriere di conoscenza nell'ERP & Robey et al. (2002) & Configurazione a KPMG, assimilazione agli utenti chiave & \hyperref[sec:9.5]{\S9.5} \\
\bottomrule
\end{xltabular}}

\section{Fonti}\label{sec:23}

\textbf{Caso e corso}
\begin{itemize}
\item Challenge KPMG – NexaCore S.p.A., materiali del corso di Comportamento Organizzativo (prof.ssa Villani), Università di Bologna.
\item Robbins, Judge e Bodega, \emph{Comportamento organizzativo}, 18ª ed., cap. 6 (McClelland).
\item Dispense del corso su cambiamento organizzativo, resistenza, ricerca-azione, struttura, cultura e motivazione.
\end{itemize}

\textbf{Trivia e motivazione}
\begin{itemize}
\item Bandura, A. (2006). Guide for Constructing Self-Efficacy Scales. In Pajares e Urdan (a cura di), \emph{Self-Efficacy Beliefs of Adolescents}, 307-337.
\item Compeau, D. R. e Higgins, C. A. (1995). Computer Self-Efficacy. \emph{MIS Quarterly}, 19(2), 189-211.
\item Schwarzer, R. e Jerusalem, M. (1995). Generalized Self-Efficacy Scale.
\item Rigotti, T., Schyns, B. e Mohr, G. (2008). A Short Version of the Occupational Self-Efficacy Scale. \emph{Journal of Career Assessment}, 16(2), 238-255.
\item Deci, E. L. e Ryan, R. M. (1985). \emph{Intrinsic Motivation and Self-Determination in Human Behavior}. Plenum.
\item Ryan, R. M. e Deci, E. L. (2000). Self-Determination Theory. \emph{American Psychologist}, 55(1), 68-78.
\item Gagné, M. et al. (2015). The Multidimensional Work Motivation Scale. \emph{European Journal of Work and Organizational Psychology}, 24(2), 178-196.
\item Adams, J. S. (1965). Inequity in Social Exchange. \emph{Advances in Experimental Social Psychology}, 2, 267-299.
\item Piderit, S. K. (2000). Rethinking Resistance and Recognizing Ambivalence. \emph{Academy of Management Review}, 25(4), 783-794.
\end{itemize}

\textbf{Valutazione}
\begin{itemize}
\item Herscovitch, L. e Meyer, J. P. (2002). Commitment to organizational change. \emph{Journal of Applied Psychology}, 87(3), 474-487.
\item HSE Management Standards Indicator Tool; metodologia INAIL per la valutazione dello stress lavoro-correlato (D.Lgs. 81/2008).
\item Topp, C. W. et al. (2015). The WHO-5 Well-Being Index: a systematic review.
\item Kristensen, T. S. et al. (2005). The Copenhagen Burnout Inventory.
\end{itemize}

\textbf{ERP, AI e automazione}
\begin{itemize}
\item Robey, D., Ross, J. W. e Boudreau, M.-C. (2002). Learning to implement enterprise systems. \emph{Journal of Management Information Systems}.
\item Casi: SACE, HolidayCo, Lidl, Accenture, IBM, TCS, Qantas (ricerche del gruppo).
\item Tempi e fasi: NBS ``Rolling out cloud ERP''; TWBS ``How long does ERP implementation take''; SAP Activate.
\item Strategie di rilascio: ERP Research ``6 ERP implementation strategies''; Panorama Consulting ``Big bang vs phased''.
\item Costi ERP: ERP Research ``Cloud ERP pricing per user 2026''; Itransition, Emerline, LeverX, Ultra Consultants sui costi di implementazione e sui costi nascosti.
\item Sicurezza: IANS/Artico ``2025 Security Budget Benchmark''; UnderDefense ``Cybersecurity budget 2026''.
\item Automazione: AIMultiple ``RPA pricing''; Costbench.
\item Shadow AI: rapporti Okta e Gartner.
\item Produttività: Hitt, Wu e Zhou (2002) sul calo dopo il go-live; Federal Reserve Bank di St. Louis (2025) sulle ore risparmiate con l'AI generativa.
\item Ripartizione 20/70/10: NOBL, ``The four types of people you'll meet when implementing change''; E. Murray (HR Reporter).
\item Costo del lavoro: Eurostat, costo orario del lavoro 2023.
\end{itemize}

\section{Nota sull'uso dell'AI}\label{sec:24}

Abbiamo usato assistenti di AI generativa per organizzare i materiali, controllare i calcoli, scrivere il codice della demo, impaginare le slide e aggregare in questo documento i testi di lavoro del gruppo. Analisi, assunzioni e scelte sono del gruppo.

\end{document}
